\documentclass[a4paper, 10pt]{article}

% --- Layout ---------------------------------------------------------------
\usepackage[margin=1in]{geometry}
\usepackage{setspace}
\usepackage{placeins}   % \FloatBarrier at every \subsection
\usepackage{float}
\usepackage{fancyhdr}

% --- Encoding & fonts -----------------------------------------------------
\usepackage[utf8]{inputenc}
\usepackage[T1]{fontenc}

% --- Math -----------------------------------------------------------------
\usepackage{amsmath}
\usepackage{amssymb}
\usepackage{mathtools}
\usepackage{bm}

% --- Tables ---------------------------------------------------------------
\usepackage{array}
\usepackage{tabularx}
\usepackage{multirow}
\usepackage{booktabs}
\usepackage{longtable}
\usepackage{colortbl}

% --- Figures --------------------------------------------------------------
\usepackage{graphicx}
\usepackage{caption}
\usepackage{subcaption}
\usepackage{xcolor}
\definecolor{invalidRow}{rgb}{1.0,0.85,0.85}

% --- Lists / boxes --------------------------------------------------------
\usepackage{enumitem}
\usepackage{tcolorbox}

% --- Hyperlinks (load last among the user-facing packages) ---------------
\usepackage[hidelinks]{hyperref}

% --- Macros / parameter handling -----------------------------------------
\usepackage{pgfkeys}

%--------------------------------------------------------------------------
% Macros
%--------------------------------------------------------------------------
\newcommand{\vect}[1]{\mathbf{#1}}
\newcommand{\vecx}[2]{\begin{bmatrix}#1\\#2\end{bmatrix}}
\newcommand{\fbplot}[5]{%
  \detokenize{Images/}#1\detokenize{_N=}#2\detokenize{,x₀=[}#3\detokenize{, }#4%
  \detokenize{],x_d=[0.0, 10.0],m=1.0,k=1.0,α=}#5\detokenize{,l₀=9.0}%
}
\definecolor{invalidRow}{rgb}{1.0,0.85,0.85}

% Convenience macros for the four scheme names — keeps notation consistent.
\newcommand{\SEoo}{SE1_1}
\newcommand{\SEot}{SE1_2}
\newcommand{\SEto}{SE2_1}
\newcommand{\SEtt}{SE2_2}

\pgfkeys{
  /algoAnalysis/.is family, /algoAnalysis,
    method/.initial=,
    k/.initial=1.0,
    h/.initial=10.0,
    N/.initial=100,
    m/.initial=1.0,
    xA/.initial=,
    xB/.initial=,
    xC/.initial=,
}

\newcommand{\plotAnalysis}[6]{%
    \edef\AlgoName{#1}%
    \edef\kval{#2}%
    \edef\hval{#3}%
    \edef\Nval{#4}%
    \edef\mval{#5}%
    \edef\x{#6}%
    \begin{figure}[H]
       \centering
       \begin{subfigure}{\textwidth}
          \centering
          \includegraphics[width=\textwidth]{Images/ControlPlot_Modified \AlgoName_k=\kval_h=\hval_N=\Nval_m=\mval_x0=\x}
          \caption{Modified \AlgoName}
       \end{subfigure}
       \vspace{0.5cm}
       \begin{subfigure}{\textwidth}
          \centering
          \includegraphics[width=\textwidth]{Images/ControlPlot_\AlgoName_k=\kval_h=\hval_N=\Nval_m=\mval_x0=\x}
          \caption{\AlgoName}
       \end{subfigure}
       \caption{Control trajectories for $\mathbf{x_0}=(\x)$ with
       $k=\kval$, $h=\hval$, $N=\Nval$, $m=\mval$.}
    \end{figure}

    \begin{figure}[H]
       \centering
       \includegraphics[width=\textwidth]{Images/Projection_Comparison_\AlgoName_k=\kval_h=\hval_N=\Nval_m=\mval_x0=\x}
       \caption{Trajectory comparison for $\mathbf{x_0}=(\x)$ with
       $k=\kval$, $h=\hval$, $N=\Nval$, $m=\mval$.}
    \end{figure}
}

\newcommand{\algoAnalysis}[1][]{%
    \pgfkeys{/algoAnalysis, xA=,xB=,xC=,#1}%
    \edef\AlgoName{\pgfkeysvalueof{/algoAnalysis/method}}%
    \edef\kval{\pgfkeysvalueof{/algoAnalysis/k}}%
    \edef\hval{\pgfkeysvalueof{/algoAnalysis/h}}%
    \edef\Nval{\pgfkeysvalueof{/algoAnalysis/N}}%
    \edef\mval{\pgfkeysvalueof{/algoAnalysis/m}}%
    \edef\xA{\pgfkeysvalueof{/algoAnalysis/xA}}%
    \edef\xB{\pgfkeysvalueof{/algoAnalysis/xB}}%
    \edef\xC{\pgfkeysvalueof{/algoAnalysis/xC}}%

    \begin{figure}[H]
       \centering
       \begin{subfigure}{\textwidth}
          \centering
          \includegraphics[width=\textwidth]{Images/IterationPlot_Modified \AlgoName_k=\kval_h=\hval_N=\Nval_m=\mval}
       \end{subfigure}
       \medskip
       \begin{subfigure}{\textwidth}
          \centering
          \includegraphics[width=\textwidth]{Images/IterationPlot_\AlgoName_k=\kval_h=\hval_N=\Nval_m=\mval}
       \end{subfigure}
       \caption{Iteration plots for \AlgoName{} and its modified variant with
       $k=\kval$, $h=\hval$, $N=\Nval$, $m=\mval$.}
       % Make the label unique per parameter set to avoid duplicate labels
       % when \algoAnalysis is invoked several times for the same method.
       \label{fig:\AlgoName-iter-k\kval-h\hval-N\Nval}
    \end{figure}

    \ifx\xA\empty\else
       \plotAnalysis{\AlgoName}{\kval}{\hval}{\Nval}{\mval}{\xA}
    \fi
    \ifx\xB\empty\else
       \plotAnalysis{\AlgoName}{\kval}{\hval}{\Nval}{\mval}{\xB}
    \fi
    \ifx\xC\empty\else
       \plotAnalysis{\AlgoName}{\kval}{\hval}{\Nval}{\mval}{\xC}
    \fi
}

% --- 4-panel grid for one quantity at fixed (α, x₀) over four Δt values --
% Args: #1 = α, #2 = x01, #3 = x02
\newcommand{\HsweepFig}[3]{%
  \begin{figure}[!htbp]
    \centering
    \begin{subfigure}{0.49\textwidth}\centering
      \includegraphics[width=\textwidth]{\fbplot{hamiltonian\_plot}{100}{#2}{#3}{#1}}
      \caption{$\Delta t = 0.1$}
    \end{subfigure}\hfill
    \begin{subfigure}{0.49\textwidth}\centering
      \includegraphics[width=\textwidth]{\fbplot{hamiltonian\_plot}{200}{#2}{#3}{#1}}
      \caption{$\Delta t = 0.05$}
    \end{subfigure}

    \medskip
    \begin{subfigure}{0.49\textwidth}\centering
      \includegraphics[width=\textwidth]{\fbplot{hamiltonian\_plot}{400}{#2}{#3}{#1}}
      \caption{$\Delta t = 0.025$}
    \end{subfigure}\hfill
    \begin{subfigure}{0.49\textwidth}\centering
      \includegraphics[width=\textwidth]{\fbplot{hamiltonian\_plot}{800}{#2}{#3}{#1}}
      \caption{$\Delta t = 0.0125$}
    \end{subfigure}
    \caption{Hamiltonian profile for $\alpha = #1$,
      $\mathbf{x_0} = (#2,\,#3)^{T}$, as the time step is refined.
      $H_i$ flattens with smaller $\Delta t$, in agreement with its
      first-integral character along the optimal trajectory.}
    \label{fig:fb-H-a#1-x#2-#3}
  \end{figure}
}

\newcommand{\UsweepFig}[3]{%
  \begin{figure}[!htbp]
    \centering
    \begin{subfigure}{0.49\textwidth}\centering
      \includegraphics[width=\textwidth]{\fbplot{control\_plot}{100}{#2}{#3}{#1}}
      \caption{$\Delta t = 0.1$}
    \end{subfigure}\hfill
    \begin{subfigure}{0.49\textwidth}\centering
      \includegraphics[width=\textwidth]{\fbplot{control\_plot}{200}{#2}{#3}{#1}}
      \caption{$\Delta t = 0.05$}
    \end{subfigure}

    \medskip
    \begin{subfigure}{0.49\textwidth}\centering
      \includegraphics[width=\textwidth]{\fbplot{control\_plot}{400}{#2}{#3}{#1}}
      \caption{$\Delta t = 0.025$}
    \end{subfigure}\hfill
    \begin{subfigure}{0.49\textwidth}\centering
      \includegraphics[width=\textwidth]{\fbplot{control\_plot}{800}{#2}{#3}{#1}}
      \caption{$\Delta t = 0.0125$}
    \end{subfigure}
    \caption{Control $u_i$ for $\alpha = #1$,
      $\mathbf{x_0} = (#2,\,#3)^{T}$, as the time step is refined.
      Boundary-layer transients sharpen and the bulk profile stabilises
      with smaller $\Delta t$.}
    \label{fig:fb-u-a#1-x#2-#3}
  \end{figure}
}

%==========================================================================
\begin{document}

\section{Introduction}
In this study, we numerically address an optimal control problem employing the Hamiltonian framework to derive the necessary conditions for optimality. The Hamiltonian formulation provides a systematic approach for deriving the governing equations.

Consider the following optimal control problem:
\begin{equation}
    \begin{aligned}
       \min \quad & J(v,u) = \int_{0}^{T} \left( \frac{v^2}{2} + \frac{\alpha}{2}u^2 \right) dt, \\
        & \dot{v} = g + \frac{u}{m} - \frac{k}{m}v, \\
       & v(0) = v_{0}.
    \end{aligned}
    \label{eq:optimal_problem}
\end{equation}
The Hamiltonian, defined as
\[
\mathcal{H}(v,u,\lambda)  = \frac{v^2}{2} + \frac{\alpha}{2}u^2 + \lambda \left( g + \frac{u}{m} - \frac{k}{m}v \right),
\]
is utilized to write the necessary optimality conditions for the above optimal control problem as follows:
\begin{equation}
    \begin{aligned}
       & \dot{v} = \frac{\partial \mathcal{H}}{\partial \lambda} (v,u,\lambda) =  g + \frac{u}{m} - \frac{k}{m}v, \\
       & \dot{\lambda} = - \frac{\partial \mathcal{H}}{\partial v} (v,u,\lambda) = -v + \frac{k}{m} \lambda, \\
       & 0 =  \frac{\partial \mathcal{H}}{\partial u} (v,u,\lambda) = \alpha u + \frac{k}{m}\lambda, \\
       & v(0) = v_{0}, \; \lambda(T) = 0.
    \end{aligned}
    \label{eq:necessary_coditions}
\end{equation}
To solve the resulting system of equations, we apply four numerical techniques, specifically the $\SEoo$, $\SEot$, $\SEto$ and $\SEtt$ methods. Each method begins with three initial points $\mathbf{z}_1 = (v_1(0),\lambda_1(0))$, $\mathbf{z}_2 =(v_2(0),\lambda_2(0))$ and $\mathbf{z}_3 = (v_3(0),\lambda_3(0))$, which form a triangle in the phase space. The area of this triangle is computed at the initial time $t = 0$, and the methods are applied, respectively, to determine the area of the triangle at the final time $t = T$. So, it allows us to identify which method effectively preserves the geometric properties of the triangle, particularly its area.

Preserving the triangle's area is important in this context, as it serves as a test for identifying symplectic algorithms. Symplectic algorithms are numerical methods designed to preserve the symplectic structure of Hamiltonian systems, which is directly related to the conservation of geometric properties like area in phase space.

By applying the $\SEoo$, $\SEot$, $\SEto$ and $\SEtt$ methods to \eqref{eq:necessary_coditions}, respectively, the resulting equations are as follows:
\begin{equation}
    \begin{aligned}
       & \SEoo = \begin{cases}
          \frac{\Delta v}{\Delta t} = \frac{v_{n+1} - v_{n}}{\Delta t}= \frac{\partial \mathcal{H}}{\partial \lambda} (v_{n+1},u_{n},\lambda_{n}) \\
          \frac{\Delta \lambda}{\Delta t} =\frac{\lambda_{n+1} - \lambda_{n}}{\Delta t} =  - \frac{\partial \mathcal{H}}{\partial v}  (v_{n+1},u_{n},\lambda_{n}) \\
          0 =  \frac{\partial \mathcal{H}}{\partial u} (v_n,u_n,\lambda_n) = \alpha u_n + \frac{k}{m}\lambda_n \\
          v(0) = v_0,\; \lambda(0) = \lambda_0.
       \end{cases} \\
       & \SEot =  \begin{cases}
            \frac{\Delta v}{\Delta t} = \frac{v_{n+1} - v_{n}}{\Delta t}= \frac{\partial \mathcal{H}}{\partial \lambda} (v_{n+1},u_{n+1},\lambda_{n}) \\
            \frac{\Delta \lambda}{\Delta t} =\frac{\lambda_{n+1} - \lambda_{n}}{\Delta t} =  - \frac{\partial \mathcal{H}}{\partial v} (v_{n+1},u_{n+1},\lambda_{n})  \\
            0 =  \frac{\partial \mathcal{H}}{\partial u} (v_n,u_n,\lambda_n) = \alpha u_n + \frac{k}{m}\lambda_n \\
            v(0) = v_0,\; \lambda(0) = \lambda_0.
       \end{cases} \\
       & \SEto = \begin{cases}
          \frac{\Delta v}{\Delta t} = \frac{v_{n+1} - v_{n}}{\Delta t}= \frac{\partial \mathcal{H}}{\partial \lambda} (v_{n},u_{n+1},\lambda_{n+1}) \\
          \frac{\Delta \lambda}{\Delta t} =\frac{\lambda_{n+1} - \lambda_{n}}{\Delta t} =  - \frac{\partial \mathcal{H}}{\partial v} (v_{n},u_{n+1},\lambda_{n+1}) \\
          0 =  \frac{\partial \mathcal{H}}{\partial u} (v_n,u_n,\lambda_n) = \alpha u_n + \frac{k}{m}\lambda_n \\
          v(0) = v_0,\; \lambda(0) = \lambda_0.
       \end{cases} \\
       & \SEtt = \begin{cases}
          \frac{\Delta v}{\Delta t} = \frac{v_{n+1} - v_{n}}{\Delta t}= \frac{\partial \mathcal{H}}{\partial \lambda} (v_{n},u_{n},\lambda_{n+1}) \\
          \frac{\Delta \lambda}{\Delta t} =\frac{\lambda_{n+1} - \lambda_{n}}{\Delta t} =  - \frac{\partial \mathcal{H}}{\partial v} (v_{n},u_{n},\lambda_{n+1}) \\
          0 =  \frac{\partial \mathcal{H}}{\partial u} (v_n,u_n,\lambda_n) = \alpha u_n + \frac{k}{m}\lambda_n \\
          v(0) = v_0,\; \lambda(0) = \lambda_0.
       \end{cases}
    \end{aligned}
    \label{eq:equation3}
\end{equation}
Before presenting numerical results, we first demonstrate that the numerical methods $\SEoo$ and $\SEto$ are symplectic integrators, which means that they exactly preserve the symplectic condition. In contrast, $\SEot$ and $\SEtt$ are not symplectic integrators.

$\bm{\SEoo:}$
\begin{align}
&   v_{n+1} = v_{n} +\Delta{t} g + \frac{\Delta t}{m} u_{n} - \frac{k\Delta t}{m}v_{n+1} \Rightarrow
\frac{m+ k \Delta t}{m} v_{n+1} = v_{n} + \Delta{t}g +  \frac{\Delta t}{m} u_{n} \Rightarrow  \nonumber \\ &
\boxed{v_{n+1} = \frac{m}{m + k\Delta{t}} v_{n}  + \frac{m\Delta{t}}{m + k\Delta{t}}g + \frac{-k\Delta{t}}{m\alpha (m+k\Delta{t})} \lambda_{n}.}  \label{eq:boxed1}\\ &
\lambda_{n+1} = \lambda_{n} - \Delta{t}v_{n+1} + \frac{k\Delta{t}}{m}\lambda_{n} \Rightarrow
\boxed{\lambda_{n+1} =\frac{m+k\Delta{t}}{m}\lambda_{n} - \Delta{t} v_{n+1}.} \label{eq:boxed2}
\end{align}
Taking differentials of both sides of equations \eqref{eq:boxed1} and \eqref{eq:boxed2} and using the linearity of the wedge product, we have
\begin{align*}
&   dv_{n+1} = \frac{m}{m + k\Delta{t}} dv_{n} + \frac{-k\Delta{t}}{m\alpha (m+k\Delta{t})} d\lambda_{n}, \\
& d\lambda_{n+1} =\frac{m+k\Delta{t}}{m} d\lambda_{n} - \Delta{t} d v_{n+1}, \\
& d\lambda_{n+1} \wedge dv_{n+1} = \frac{m + k\Delta{t}}{m} d\lambda_{n} \wedge dv_{n+1} =  \frac{m+k\Delta{t}}{m} d\lambda_{n}  \wedge \left(
 \frac{m}{m + k\Delta{t}} dv_{n} + \frac{-k\Delta{t}}{m\alpha (m+k\Delta{t})} d\lambda_{n} \right) \Rightarrow   \\
& \boxed{ d\lambda_{n+1} \wedge dv_{n+1} =  \frac{m+k\Delta{t}}{m} d\lambda_{n}  \wedge  \frac{m}{m + k\Delta{t}} dv_{n}
 =  d\lambda_{n} \wedge dv_{n}}.
\end{align*}

$\bm{\SEto:}$
\begin{align}
&   v_{n+1} = \frac{m-k\Delta{t}}{m} v_{n} + \Delta{t}g + \frac{\Delta{t}}{m} u_{n+1} \Rightarrow \nonumber \\ &
    \boxed{v_{n+1} = \frac{m-k\Delta{t}}{m} v_{n}  + \Delta{t}g + \frac{-k\Delta{t}}{m^2\alpha} \lambda_{n+1}.}  \label{eq:boxed3}\\ &
    \boxed{\lambda_{n+1} = \frac{m}{m-\Delta{t}k}(-\Delta{t})v_{n} + \frac{m}{m-\Delta{t}k}\lambda_{n}.}  \label{eq:boxed4}
\end{align}
Similar to $\SEoo$, by taking differentials of both sides of equations \eqref{eq:boxed3} and \eqref{eq:boxed4} and using the wedge product, we have
\begin{align*}
& dv_{n+1} = \frac{m-k\Delta{t}}{m} dv_{n} + \frac{-k\Delta{t}}{m^2\alpha} d\lambda_{n+1}, \\
& d\lambda_{n+1} = \frac{m}{m-\Delta{t}k}(-\Delta{t}) dv_{n} + \frac{m}{m-\Delta{t}k} d\lambda_{n}, \\
& \boxed{dv_{n+1} \wedge d\lambda_{n+1} = \frac{m-k\Delta{t}}{m} dv_{n} \wedge d\lambda_{n+1} = dv_{n} \wedge d\lambda_{n}.}
\end{align*}

$\bm{\SEot:}$
\begin{align}
&   v_{n+1} = \frac{m}{m + k\Delta{t}} v_{n} +  \frac{m}{m + k\Delta{t}}\Delta{t}g + \frac{\Delta{t}}{m + k\Delta{t}} u_{n+1} \Rightarrow  \nonumber \\ &
\boxed{v_{n+1}  = \frac{m}{m + k\Delta{t}} v_{n} + \frac{m}{m + k\Delta{t}}\Delta{t}g + \frac{-k\Delta{t}}{m\alpha(m + k\Delta{t})} \lambda_{n+1}.} \label{eq:boxed5} \\ &
\boxed{\lambda_{n+1} = \frac{m+k\Delta{t}}{m} \lambda_{n} + (-\Delta{t})v_{n+1}.} \label{eq:boxed6}
\end{align}
Now by taking differentials of both sides of equations \eqref{eq:boxed5} and \eqref{eq:boxed6}, we have
\begin{align}
    &  dv_{n+1}  = \frac{m}{m + k\Delta{t}} dv_{n}  + \frac{-k\Delta{t}}{m\alpha(m + k\Delta{t})} d\lambda_{n+1} \label{eq:boxed7}, \\
    & d\lambda_{n+1} = \frac{m+k\Delta{t}}{m} d\lambda_{n} + (-\Delta{t}) dv_{n+1}. \label{eq:boxed8}
\end{align}
From equations \eqref{eq:boxed7} and \eqref{eq:boxed8}, $dv_{n+1}$ and $d\lambda_{n+1}$ can be expressed in terms of $dv_{n}$ and $d\lambda_{n}$, so we have:
\begin{align}
    &  dv_{n+1}  = \frac{m^2\alpha}{m\alpha(m + k\Delta{t})-k(\Delta{t})^2} dv_{n}  +  \frac{-k\Delta{t}(m+k\Delta{t})}{m^2\alpha(m+k\Delta{t})-km(\Delta{t})^2}d\lambda_{n} \label{eq_9} \\
    & d\lambda_{n+1} = \frac{\alpha(m+k\Delta{t})^2}{m\alpha(m+k\Delta{t})-k(\Delta{t})^2}d\lambda_{n} + \frac{-m^2\alpha\Delta{t}}{m\alpha(m+k\Delta{t})-k(\Delta{t})^2} dv_{n} \label{eq_10}
\end{align}
If we take the wedge product of \eqref{eq_9} with \eqref{eq_10}, we get
\begin{align}
    \boxed{dv_{n+1} \wedge d\lambda_{n+1} = \left(1 + \frac{k(\Delta{t})^2}{m\alpha(m+k\Delta{t})-k(\Delta{t})^2}\right) dv_{n} \wedge d\lambda_{n}. }
\end{align}

$\bm{\SEtt:}$
\begin{align}
  & v_{n+1} = g\Delta{t} + \tfrac{m-k\Delta{t}}{m} v_{n} + \tfrac{-k\Delta{t}}{m^2\alpha}\lambda_{n}, \label{eq:se2_2_11} \\
  & \lambda_{n+1} = \tfrac{-m\Delta{t}}{m-k\Delta{t}}v_{n} + \tfrac{m}{m-k\Delta{t}}\lambda_{n}. \label{eq:se2_2_12}
\end{align}
Taking differentials of both sides of equations \eqref{eq:se2_2_11} and \eqref{eq:se2_2_12} and using the wedge product, we have
\begin{equation}
    dv_{n+1} \wedge d\lambda_{n+1} = \left(1 - \frac{k(\Delta{t})^2}{m\alpha(m-k\Delta{t})}\right) dv_{n} \wedge d\lambda_{n}.
    \label{eq:equation4}
\end{equation}

\subsection{Linear optimal control problem}
Let us consider the following linear optimal control problem:
\begin{equation}
    \label{eq_18}
\begin{aligned}
    J(\bm{v},\bm{u}) & = \frac{1}{2} \int_{a}^b \bm{v}(t)^{T} \bm{Q} \bm{v}(t) + \bm{u}(t)^{T} \bm{R} \bm{u}(t)\; dt \\
    & \dot{\bm{v}}(t) = \bm{A} \bm{v}(t) + \bm{B}\bm{u}(t), \\
    & \bm{v}(a) = \bm{v}_{0},
\end{aligned}
\end{equation}
where $\bm{Q}$ is a real symmetric positive semi-definite matrix and $\bm{R}$ is a real symmetric positive definite matrix.

The Hamiltonian is
\begin{equation}\label{eq_19}
    \mathcal{H}(\bm{v},\bm{u},\bm{\lambda}) =  \frac{1}{2} \bm{v}(t)^{T} \bm{Q} \bm{v}(t) + \frac{1}{2} \bm{u}(t)^{T} \bm{R} \bm{u}(t) + \bm{\lambda}^{T} \bm{A}\bm{v}(t) + \bm{\lambda}^{T} \bm{B}\bm{u}(t),
\end{equation}
and necessary conditions for optimality are
\begin{align}
    & \dot{\bm{v}}(t) = \frac{\partial \mathcal{H}(\bm{v},\bm{u},\bm{\lambda})}{\partial\bm{\lambda}}  = \bm{A} \bm{v}(t) + \bm{B}\bm{u}(t), \label{eq_20} \\
    & \dot{\bm{\lambda}}(t) =  -\frac{\partial \mathcal{H}(\bm{v},\bm{u},\bm{\lambda})}{\partial \bm{v}} = -\bm{Q}\bm{v}(t) - \bm{A}^{T}\bm{\lambda}(t), \label{eq_21} \\
    & \bm{0} =  \frac{\partial \mathcal{H}(\bm{v},\bm{u},\bm{\lambda})}{\partial \bm{u}} = \bm{R}\bm{u}(t) + \bm{B}^{T}\bm{\lambda}(t). \label{eq_22}
\end{align}
Equation \eqref{eq_22} can be solved for $\bm{u}$ to give
\begin{align}
    \bm{u}(t) = -\bm{R}^{-1}\bm{B}^{T}\bm{\lambda}(t),
\end{align}
the existence of $\bm{R}^{-1}$ is assured, since $\bm{R}$ is a positive definite matrix.

Let us apply the $\SEoo$ and $\SEto$ schemes to equations \eqref{eq_20}--\eqref{eq_22}, respectively. We want to show that both are symplectic integrators, that is, they preserve the symplectic condition, i.e.\ $d\bm{v}_{n+1} \wedge d\bm{\lambda}_{n+1} = d\bm{v}_{n} \wedge d\bm{\lambda}_{n}$ or $\det\!\left(\frac{\partial(\bm{v}_{n+1},\bm{\lambda}_{n+1})}{\partial(\bm{v}_{n},\bm{\lambda}_{n})}\right) = 1.$

$\bm{\SEoo:}$
\begin{align}
&   \frac{\bm{v}_{n+1} - \bm{v}_n}{\Delta{t}} = \bm{A}\bm{v}_{n+1} + \bm{B}\bm{u}_{n} =  \bm{A}\bm{v}_{n+1} - \bm{B}\bm{R}^{-1}\bm{B}^{T}\bm{\lambda}_{n} \Rightarrow \nonumber \\ &
    \bm{v}_{n+1} = (\bm{I} - \Delta{t}\bm{A})^{-1} \bm{v}_{n} - \Delta{t}(\bm{I} - \Delta{t}\bm{A})^{-1}\bm{B}\bm{R}^{-1}\bm{B}^{T}\bm{\lambda}_{n}, \label{eq_24} \\ &
       \frac{\bm{\lambda}_{n+1} - \bm{\lambda}_n}{\Delta{t}} = -\bm{Q}\bm{v}_{n+1} - \bm{A}^{T}\bm{\lambda}_{n} \Rightarrow \nonumber \\ &
       \bm{\lambda}_{n+1} = -\Delta{t}\bm{Q}\bm{v}_{n+1} +(\bm{I} - \Delta{t}\bm{A}^{T})\bm{\lambda}_{n}. \label{eq_25}
\end{align}
Taking differentials of \eqref{eq_24} and \eqref{eq_25}, respectively, we obtain the following implicit system of equations:
\begin{align}
    &  d\bm{v}_{n+1} = (\bm{I} - \Delta{t}\bm{A})^{-1} d\bm{v}_{n} - \Delta{t}(\bm{I} - \Delta{t}\bm{A})^{-1}\bm{B}\bm{R}^{-1}\bm{B}^{T}d\bm{\lambda}_{n}, \label{eq_26} \\
    & d\bm{\lambda}_{n+1} = -\Delta{t}\bm{Q}d\bm{v}_{n+1} +(\bm{I} - \Delta{t}\bm{A}^{T})d\bm{\lambda}_{n}. \label{eq_27}
\end{align}
By taking the wedge product of \eqref{eq_26} with \eqref{eq_27} from the right, we have
\begin{align}
       d\bm{v}_{n+1} \wedge   d\bm{\lambda}_{n+1}  = &  (\bm{I} - \Delta{t}\bm{A})^{-1} d\bm{v}_{n} \wedge  \left(-\Delta{t}\bm{Q}d\bm{v}_{n+1}\right) +   (\bm{I} - \Delta{t}\bm{A})^{-1} d\bm{v}_{n} \wedge (\bm{I} - \Delta{t}\bm{A}^{T})d\bm{\lambda}_{n}  +  \nonumber \\ &
          \left(-\Delta{t}(\bm{I} - \Delta{t}\bm{A})^{-1}\bm{B}\bm{R}^{-1}\bm{B}^{T}d\bm{\lambda}_{n}\right) \wedge \left(-\Delta{t}\bm{Q}d\bm{v}_{n+1} \right). \label{eq_28}
\end{align}
Upon plugging \eqref{eq_26} into \eqref{eq_28} and using $d\bm{a} \wedge d(\bm{M}d\bm{b}) = (\bm{M}^{T}d\bm{a}) \wedge d\bm{b}$, we obtain the following equality
\begin{align}
    d\bm{v}_{n+1} \wedge d\bm{\lambda}_{n+1} = &  (\bm{I} - \Delta{t}\bm{A})^{-1} d\bm{v}_{n} \wedge \Delta{t}^2\bm{Q}(\bm{I} - \Delta{t}\bm{A})^{-1}\bm{B}\bm{R}^{-1}\bm{B}^{T}d\bm{\lambda}_{n} + d\bm{v}_{n} \wedge d\bm{\lambda}_{n} + \nonumber \\
    & - \Delta{t}(\bm{I} - \Delta{t}\bm{A})^{-1}\bm{B}\bm{R}^{-1}\bm{B}^{T}d\bm{\lambda}_{n} \wedge -\Delta{t}\bm{Q}(\bm{I} - \Delta{t}\bm{A})^{-1} d\bm{v}_n   \nonumber \\
    & =  d\bm{v}_{n} \wedge d\bm{\lambda}_{n} - \Delta{t}^{2}(\bm{I}- \Delta{t}\bm{A})^{-T}\bm{Q}(\bm{I} - \Delta{t}\bm{A})^{-1}\bm{B}\bm{R}^{-1}\bm{B}^{T}d\bm{\lambda}_{n} \wedge d\bm{v}_n + \nonumber \\
    & \Delta{t}^{2}(\bm{I}- \Delta{t}\bm{A})^{-T}\bm{Q}^{T}(\bm{I} - \Delta{t}\bm{A})^{-1}\bm{B}\bm{R}^{-1}\bm{B}^{T}d\bm{\lambda}_{n} \wedge d\bm{v}_n  \nonumber \\
    & =  d\bm{v}_n \wedge d\bm{\lambda}_{n}.
\end{align}

$\bm{\SEto:}$
\begin{align}
    &  \bm{v}_{n+1} = (\bm{I} + \Delta{t}\bm{A}) \bm{v}_{n} - \Delta{t}\bm{B}\bm{R}^{-1}\bm{B}^{T}\bm{\lambda}_{n+1}, \label{eq_30} \\
    & \bm{\lambda}_{n+1} = -\Delta{t}(\bm{I}+\Delta{t}\bm{A}^{T})^{-1}\bm{Q}\bm{v}_{n} +(\bm{I} + \Delta{t}\bm{A}^{T})^{-1}\bm{\lambda}_{n}, \label{eq_31} \\
    & d\bm{v}_{n+1} = (\bm{I} + \Delta{t}\bm{A}) d\bm{v}_{n} - \Delta{t}\bm{B}\bm{R}^{-1}\bm{B}^{T}d\bm{\lambda}_{n+1}, \label{eq_32} \\
    & d\bm{\lambda}_{n+1} = -\Delta{t}(\bm{I}+\Delta{t}\bm{A}^{T})^{-1}\bm{Q}d\bm{v}_{n} +(\bm{I} + \Delta{t}\bm{A}^{T})^{-1}d\bm{\lambda}_{n}. \label{eq_33}
\end{align}
By taking the wedge product of \eqref{eq_32} with \eqref{eq_33} we obtain
\begin{align}
    d\bm{\lambda}_{n+1} \wedge     d\bm{v}_{n+1} = &  \left(-\Delta{t}(\bm{I}+\Delta{t}\bm{A}^{T})^{-1}\bm{Q}d\bm{v}_{n}\right) \wedge  \left(- \Delta{t}\bm{B}\bm{R}^{-1}\bm{B}^{T}d\bm{\lambda}_{n+1} \right) + \nonumber \\
    & (\bm{I} + \Delta{t}\bm{A}^{T})^{-1}d\bm{\lambda}_{n} \wedge (\bm{I} + \Delta{t}\bm{A})d\bm{v}_{n} + \nonumber \\
    & (\bm{I} + \Delta{t}\bm{A}^{T})^{-1}d\bm{\lambda}_{n} \wedge \left(-\Delta{t}\bm{B}\bm{R}^{-1}\bm{B}^{T}d\bm{\lambda}_{n+1}\right). \label{eq_34}
\end{align}
Upon plugging \eqref{eq_33} into \eqref{eq_34}, we obtain the desired equality \eqref{eq_35}.
\begin{align}
       d\bm{\lambda}_{n+1} \wedge     d\bm{v}_{n+1}  = & d\bm{\lambda}_{n} \wedge d\bm{v}_{n} + \nonumber \\
       & -\Delta{t}^2\left((\bm{I} + \Delta{t}\bm{A}^{T})^{-1}\bm{Q}\right)^{T}\bm{B}\bm{R}^{-1}\bm{B}^{T}(\bm{I} + \Delta{t}\bm{A}^{T})^{-1} d\bm{\lambda}_{n} \wedge d\bm{v}_{n}  + \nonumber \\
       & \Delta{t}^2 \left(\bm{B}\bm{R}^{-1}\bm{B}^{T}(\bm{I} + \Delta{t}\bm{A}^{T})^{-1}\bm{Q}\right)^{T} (\bm{I} + \Delta{t}\bm{A}^{T})^{-1} d\bm{\lambda}_{n} \wedge d\bm{v}_{n}  = \nonumber \\
       & d\bm{\lambda}_{n} \wedge d\bm{v}_{n}. \label{eq_35}
\end{align}

\section{Numerical Results (linear case)}

In this section, we present numerical results (summarized in Table~\ref{tab:area-linear}) that compare the initial and final areas computed by each method. Figures~\ref{fig:se11}--\ref{fig:se22} depict the area of the triangles calculated in discrete time steps $(t = i\Delta{t},\; \Delta{t} = 0.1)$ with different $\alpha$ over the interval $[0, 10]$. This stepwise analysis highlights how the area varies in the time range for each approach.

\begin{table}[ht!]
\centering
\begin{tabular}{|c|c|c|c|}
   \hline
Method & $\alpha$ & $A_0$ & $A_T$ \\
   \hline
   \multirow{6}{3em}{$\SEoo$}
   & 10  & 0.1000  & 0.1000 \\
   & 7   & 0.1000  & 0.1000 \\
   & 5   & 0.1000  & 0.1000 \\
   & 2   & 0.1000  & 0.1000 \\
   & 1   & 0.1000  & 0.1000 \\
   & 0.5 & 0.1000  & 0.0991 \\
   \hline
   \multirow{6}{3em}{$\SEot$}
   & 10  & 0.1000  & 0.1095 \\
   & 7   & 0.1000  & 0.1139 \\
   & 5   & 0.1000  & 0.1200 \\
   & 2   & 0.1000  & 0.1577 \\
   & 1   & 0.1000  & 0.2493 \\
   & 0.5 & 0.1000  & 0.4609 \\
   \hline
   \multirow{6}{3em}{$\SEto$}
   & 10  & 0.1000  & 0.1000 \\
   & 7   & 0.1000  & 0.1000 \\
   & 5   & 0.1000  & 0.1000 \\
   & 2   & 0.1000  & 0.1000 \\
   & 1   & 0.1000  & 0.1001 \\
   & 0.5 & 0.1000  & 0.2500 \\
   \hline
   \multirow{6}{3em}{$\SEtt$}
   & 10  & 0.1000  & 0.0895 \\
   & 7   & 0.1000  & 0.0853 \\
   & 5   & 0.1000  & 0.0801 \\
   & 2   & 0.1000  & 0.0573 \\
   & 1   & 0.1000  & 0.0327 \\
   & 0.5 & 0.1000  & 0.0117 \\
   \hline
\end{tabular}
\caption{Application of methods $\SEoo$, $\SEot$, $\SEto$ and $\SEtt$ to system \eqref{eq:necessary_coditions} on the interval $[0,10]$ with three initial points $\mathbf{z}_1 = (10,0.1)^{T}$, $\mathbf{z}_2 =(11,0.1)^{T}$ and $\mathbf{z}_3 = (10,-0.1)^{T}$; in addition, $k =1$, $m =1$, $\Delta t =0.1$, $T = 10$. $A_0$ and $A_T$ denote the areas of the triangle at time $t = 0$ and $t = T = 10$, respectively.}
\label{tab:area-linear}
\end{table}

\begin{figure}[ht]
    \centering
    \includegraphics[width=\textwidth]{Images/se11.pdf}
    \caption{$\SEoo$ method. Triangle area vs.\ time, computed at discrete steps, $t = i\Delta{t},\; \Delta{t} = 0.1$, over interval $[0,10]$.}
    \label{fig:se11}
\end{figure}

\begin{figure}[ht]
    \centering
    \includegraphics[width=\textwidth]{Images/se12.pdf}
    \caption{$\SEot$ method. Triangle area vs.\ time, computed at discrete steps, $t = i\Delta{t},\; \Delta{t} = 0.1$, over interval $[0,10]$.}
    \label{fig:se12}
\end{figure}

\begin{figure}[ht]
    \centering
    \includegraphics[width=\textwidth]{Images/se21.pdf}
    \caption{$\SEto$ method. Triangle area vs.\ time, computed at discrete steps, $t = i\Delta{t},\; \Delta{t} = 0.1$, over interval $[0,10]$.}
    \label{fig:se21}
\end{figure}

\begin{figure}[ht]
    \centering
    \includegraphics[width=\textwidth]{Images/se22.pdf}
    \caption{$\SEtt$ method. Triangle area vs.\ time, computed at discrete steps, $t = i\Delta{t},\; \Delta{t} = 0.1$, over interval $[0,10]$.}
    \label{fig:se22}
\end{figure}

\section{Optimal control problem}

Consider the following basic optimal control problem,
\begin{equation}
    \begin{aligned}
       & \min \int_{a}^b\, j(\bm{y},\bm{u})\, dt \\
       & \dot{\bm{y}} = \bm{f}(\bm{y},\bm{u}) \\
       & \bm{y}(0) = \bm{y}_{0}.
    \end{aligned}
    \label{eq:equation5}
\end{equation}

The necessary conditions are derived from the Hamiltonian $H$, defined as
\begin{align}
    H(\bm{y},\bm{u},\bm{\lambda}) = j(\bm{y},\bm{u}) + \bm{\lambda}^{T}\bm{f}(\bm{y},\bm{u}).
\end{align}
The necessary optimality conditions can be written in terms of the Hamiltonian:
\begin{align}
    & \dot{\bm{y}} = H_{\bm{\lambda}}  = \frac{\partial H}{\partial \bm{\lambda}}(\bm{y},\bm{u},\bm{\lambda})\\
    & \dot{\bm{\lambda}} = -H_{\bm{y}} = - \frac{\partial H}{\partial \bm{y}}(\bm{y},\bm{u},\bm{\lambda})  \\
    & \bm{0} = H_{\bm{u}}= \frac{\partial H}{\partial \bm{u}}(\bm{y},\bm{u},\bm{\lambda}). \label{eq_40}
\end{align}
We want to show that $\SEoo$ and $\SEto$ schemes are symplectic integrators, that is, $d\bm{y}_{n+1} \wedge d\bm{\lambda}_{n+1} = d\bm{y}_{n} \wedge d\bm{\lambda}_{n}$ if $H_{\bm{yu}} = \frac{\partial^{2}H}{\partial \bm{u}\partial\bm{y}} = \bm{0}$.

$\bm{\SEoo:}$
Let $\bm{u}$ be a function of $\bm{y}$ and $\bm{\lambda}$, i.e.\ $\bm{u} = \bm{g}(\bm{y},\bm{\lambda})$, so from \eqref{eq_40} we have
\begin{align}
    & H_{\bm{uy}} + H_{\bm{uu}}\bm{g_y} = 0 \Rightarrow \bm{g_y} = - H_{\bm{uu}}^{-1}H_{\bm{uy}}  \label{eq_41} \\
    & H_{\bm{u\lambda}} + H_{\bm{uu}}\bm{g_{\lambda}} = 0 \Rightarrow \bm{g_{\lambda}} = - H_{\bm{uu}}^{-1}H_{\bm{u\lambda}}, \label{eq_42} \\
    & d\bm{u} =  - H_{\bm{uu}}^{-1}H_{\bm{uy}} d\bm{y}  - H_{\bm{uu}}^{-1}H_{\bm{u\lambda}} d\bm{\lambda} \label{eq_43}
\end{align}
where $\bm{g_y} = \frac{\partial\bm{g}}{\partial\bm{y}}$ and $\bm{g_{\lambda}} = \frac{\partial\bm{g}}{\partial\bm{\lambda}}$, respectively.
\begin{align}
     \bm{y}_{n+1} = \bm{y}_n + \Delta{t} H_{\bm{\lambda}}(\bm{y}_{n+1},\bm{u}_n,\bm{\lambda}_n). \label{eq_44}
\end{align}
Taking the differential of \eqref{eq_44} and plugging \eqref{eq_41}--\eqref{eq_43} into \eqref{eq_45} we obtain the equation:
\begin{align}
    d\bm{y}_{n+1} = d\bm{y}_n + \Delta{t}(H_{\bm{\lambda y}}d\bm{y}_{n+1} + H_{\bm{\lambda u}}d\bm{u}_{n} + H_{\bm{\lambda \lambda}}d\bm{\lambda}_{n}), \label{eq_45}
\end{align}
\begin{align}
    &  d\bm{y}_{n+1} = d\bm{y}_n + \Delta{t}H_{\bm{\lambda y}}d\bm{y}_{n+1} -  \Delta{t} H_{\bm{\lambda u}}H_{\bm{uu}}^{-1}H_{\bm{uy}}d\bm{y}_n - \Delta{t} H_{\bm{\lambda u}} H_{\bm{uu}}^{-1}H_{\bm{u\lambda}}d\bm{\lambda}_n + \Delta{t} H_{\bm{\lambda \lambda}}d\bm{\lambda}_{n} \\
    & \boxed{d\bm{y}_{n+1} = (\bm{I} - \Delta{t}H_{\bm{\lambda y}})^{-1}(\bm{I} - \Delta{t}\bm{A}) d\bm{y}_n +\Delta{t}(\bm{I} - \Delta{t}H_{\bm{\lambda y}})^{-1}(-\bm{B} + H_{\bm{\lambda \lambda}})d\bm{\lambda}_{n}}, \label{eq_47}
\end{align}
where $\bm{A} = H_{\bm{\lambda u}}H_{\bm{uu}}^{-1}H_{\bm{uy}}$ and $\bm{B} = H_{\bm{\lambda u}} H_{\bm{uu}}^{-1}H_{\bm{u\lambda}}$.
\begin{align}
    & \bm{\lambda}_{n+1} = \bm{\lambda}_n - \Delta{t} H_{\bm{y}}(\bm{y}_{n+1},\bm{u}_n,\bm{\lambda}_n) \\
    & d\bm{\lambda}_{n+1} = d\bm{\lambda}_n - \Delta{t}(H_{\bm{yy}}d\bm{y}_{n+1} + H_{\bm{y u}}d\bm{u}_{n} + H_{\bm{y \lambda}}d\bm{\lambda}_{n}) \label{eq_49}
\end{align}
Upon plugging \eqref{eq_41}--\eqref{eq_43} into \eqref{eq_49}, we have
\begin{align}
    &  d\bm{\lambda}_{n+1} = d\bm{\lambda}_n - \Delta{t}H_{\bm{yy}}d\bm{y}_{n+1} + \Delta{t} H_{\bm{y u}}H_{\bm{uu}}^{-1}H_{\bm{uy}}d\bm{y}_{n} + \Delta{t} H_{\bm{y u}}H_{\bm{uu}}^{-1}H_{\bm{u\lambda}}d \bm{\lambda}_{n} -\Delta{t}  H_{\bm{y \lambda}}d\bm{\lambda}_{n} \\
    & d\bm{\lambda}_{n+1} = (\bm{I}-\Delta{t}H_{\bm{y\lambda}} + \Delta{t}\bm{A}^{T})d\bm{\lambda}_{n} + \Delta{t}\bm{D} d\bm{y}_n - \Delta{t}H_{\bm{yy}}d\bm{y}_{n+1} \label{eq_51}
\end{align}
where $\bm{D} = H_{\bm{yu}}H^{-1}_{\bm{uu}}H_{\bm{uy}}$.

By taking the wedge product of \eqref{eq_51} with $d\bm{y}_{n+1}$ we obtain
\begin{align*}
    d\bm{\lambda}_{n+1} \wedge d\bm{y}_{n+1} = & (\bm{I} - \Delta{t}H_{\bm{y\lambda}} + \Delta{t}\bm{A}^{T})d\bm{\lambda}_{n}\wedge (\bm{I}- \Delta{t}H_{\bm{\lambda y}})^{-1}(\bm{I} - \Delta{t}\bm{A})d\bm{y}_{n} + \\
    & \Delta{t}(\bm{I}- \Delta{t}H_{\bm{\lambda y}})^{-1}(\bm{B} - H_{\bm{\lambda \lambda}}) d\bm{\lambda}_{n} \wedge \Delta{t}\bm{D} d\bm{y}_{n}.
\end{align*}
If $H_{\bm{yu}} = \bm{0}$, then $\bm{A} = \bm{0}$, $\bm{A}^{T} = \bm{0}$, $\bm{D} = \bm{0}$ so we have
\begin{align}
    d\bm{\lambda}_{n+1} \wedge d\bm{y}_{n+1} = & (\bm{I} - \Delta{t}H_{\bm{y\lambda}})d\bm{\lambda}_{n}\wedge (\bm{I}- \Delta{t}H_{\bm{\lambda y}})^{-1}d\bm{y}_{n}  =  \\
    & (\bm{I}- \Delta{t}H_{\bm{\lambda y}})^{-T}(\bm{I} - \Delta{t}H_{\bm{y\lambda}}) d\bm{\lambda}_{n} \wedge d\bm{y}_{n} = \\
    & (\underbrace{\bm{I}- \Delta{t}H_{\bm{\lambda y}}^{T}}_{H_{\bm{y\lambda}} = H_{\bm{\lambda y}}^{T}})^{-1}(\bm{I} - \Delta{t}H_{\bm{y\lambda}}) d\bm{\lambda}_{n} \wedge d\bm{y}_{n}= \\
    & (\bm{I}- \Delta{t}H_{\bm{y \lambda}})^{-1}(\bm{I} - \Delta{t}H_{\bm{y\lambda}}) d\bm{\lambda}_{n} \wedge d\bm{y}_{n} \\
    & =  d\bm{\lambda}_{n} \wedge d\bm{y}_{n}.
\end{align}

$\bm{\SEto:}$
\begin{align}
      \bm{y}_{n+1} = \bm{y}_n + \Delta{t} H_{\bm{\lambda}}(\bm{y}_{n},\bm{u}_{n+1},\bm{\lambda}_{n+1})
\end{align}
\begin{equation}\label{eq_58}
\begin{aligned}
    d\bm{y}_{n+1} =  & d\bm{y}_n + \Delta{t}H_{\bm{\lambda y}}d\bm{y}_{n} -  \Delta{t} H_{\bm{\lambda u}}H_{\bm{uu}}^{-1}H_{\bm{uy}}d\bm{y}_{n+1} -  \\
    & \Delta{t} H_{\bm{\lambda u}} H_{\bm{uu}}^{-1}H_{\bm{u\lambda}}d\bm{\lambda}_{n+1} + \Delta{t} H_{\bm{\lambda \lambda}}d\bm{\lambda}_{n+1} \\
    & = (\bm{I} + \Delta{t}H_{\bm{\lambda u}}H^{-1}_{\bm{uu}}H_{\bm{uy}})^{-1}(\bm{I} + \Delta{t}H_{\bm{\lambda y}}) d\bm{y}_n + \\
    & (\bm{I} + \Delta{t}H_{\bm{\lambda u}}H^{-1}_{\bm{uu}}H_{\bm{uy}})^{-1}(-\Delta{t}H_{\bm{\lambda u}}H_{\bm{uu}}H_{\bm{u\lambda}} + \Delta{t}H_{\bm{\lambda\lambda}})  d\bm{\lambda}_{n+1},
\end{aligned}
\end{equation}
\begin{align}
      \bm{\lambda}_{n+1} = \bm{\lambda}_n - \Delta{t} H_{\bm{y}}(\bm{y}_{n},\bm{u}_{n+1},\bm{\lambda}_{n+1})
\end{align}
\begin{equation}\label{eq_60}
   \begin{aligned}
         d\bm{\lambda}_{n+1} =  & d\bm{\lambda}_n -\Delta{t}H_{\bm{y y}}d\bm{y}_{n} + \Delta{t} H_{\bm{yu}}H_{\bm{uu}}^{-1}H_{\bm{uy}}d\bm{y}_{n+1} +  \\
      & \Delta{t} H_{\bm{yu}}H_{\bm{uu}}^{-1}H_{\bm{u\lambda}}d\bm{\lambda}_{n+1} - \Delta{t} H_{\bm{y \lambda}}d\bm{\lambda}_{n+1} \\
      & =  (\bm{I} - \Delta{t}H_{\bm{yu}}H^{-1}_{\bm{uu}}H_{\bm{u\lambda}}+\Delta{t}H_{\bm{y\lambda}})^{-1} d\bm{\lambda}_n - \\
      & \Delta{t} (\bm{I} - \Delta{t}H_{\bm{y u}}H^{-1}_{\bm{uu}}H_{\bm{u\lambda}}+\Delta{t}H_{\bm{y\lambda}})^{-1}H_{\bm{yy}}d\bm{y}_{n} + \\
      & \Delta{t}(\bm{I} - \Delta{t}H_{\bm{y u}}H^{-1}_{\bm{uu}}H_{\bm{u\lambda}}+\Delta{t}H_{\bm{y\lambda}})^{-1}H_{\bm{yu}}H^{-1}_{\bm{uu}}H_{\bm{uy}}d\bm{y}_{n+1}.
   \end{aligned}
\end{equation}

If $H_{\bm{uy}} = \bm{0}$, the equations \eqref{eq_58} and \eqref{eq_60} simplify to
\begin{align}
      & d\bm{y}_{n+1} = (\bm{I} + \Delta{t}H_{\bm{\lambda y}}) d\bm{y}_{n} + (-\Delta{t}H_{\bm{\lambda u}}H^{-1}_{\bm{uu}}H_{\bm{u\lambda}} + \Delta{t} H_{\bm{\lambda\lambda}}) d\bm{\lambda}_{n+1} \label{eq_61}  \\
      & d\bm{\lambda}_{n+1} = (\bm{I} + \Delta{t}H_{\bm{y\lambda}})^{-1} d\bm{y}_{n} -\Delta{t}(\bm{I} + \Delta{t}H_{\bm{y\lambda}})^{-1}H_{\bm{yy}}d\bm{y}_n \label{eq_62}
\end{align}
Taking the wedge product of \eqref{eq_61} with $d\bm{\lambda}_{n+1}$ and substituting equation \eqref{eq_62} into the resulting expression, we obtain:
\begin{equation}
\begin{aligned}
    d\bm{y}_{n+1} \wedge d\bm{\lambda}_{n+1} = &  (\bm{I} + \Delta{t}H_{\bm{\lambda y}}) d\bm{y}_{n} \wedge  (\bm{I} + \Delta{t}H_{\bm{y\lambda}})^{-1} d\bm{\lambda}_{n}   \\
    & = (\underbrace{\bm{I} + \Delta{t}H_{\bm{y\lambda}}^{T}}_{H_{\bm{y\lambda}}^{T}=H_{\bm{\lambda y}}})^{-1} (\bm{I} + \Delta{t}H_{\bm{\lambda y}}) d\bm{y}_{n} \wedge d\bm{\lambda}_n \\
    & = d\bm{y}_n \wedge d\bm{\lambda}_n.
\end{aligned}
\label{eq:equation6}
\end{equation}

\textbf{Modified MidPoint:} $(\bm{M.M.P})$
\begin{align}
    & \bm{y}_{n+1} = \bm{y}_{n} + \Delta{t}H_{\bm{\lambda}}(\bm{y_{n+\frac{1}{2}}}, \bm{u}(\bm{y_{n+\frac{1}{2}}},\bm{\lambda_{n+\frac{1}{2}}}),\bm{\lambda_{n+\frac{1}{2}}}), \nonumber \\
    & d\bm{y}_{n+1} = d\bm{y}_{n} + \tfrac{\Delta{t}}{2}H_{\bm{\lambda}\bm{y}}(d\bm{y}_{n+1}+d\bm{y}_{n})+ \tfrac{\Delta{t}}{2}H_{\bm{\lambda}\bm{u}}d\bm{u}+ \underbrace{\tfrac{\Delta{t}}{2}H_{\bm{\lambda}\bm{\lambda}}(d\bm{\lambda}_{n+1}+d\bm{\lambda}_{n})}_{\bm{0}}, \nonumber \\
    & d\bm{y}_{n+1} = d\bm{y}_{n}  +  \tfrac{\Delta{t}}{2}H_{\bm{\lambda}\bm{y}}(d\bm{y}_{n+1}+d\bm{y}_{n})
    + \tfrac{\Delta{t}}{2}H_{\bm{\lambda}\bm{u}}\!\left( -H_{\bm{uu}}^{-1}H_{\bm{uy}}(d\bm{y}_{n+1}+d\bm{y}_{n})-H_{\bm{uu}}^{-1}H_{\bm{u\lambda}}(d\bm{\lambda}_{n+1}+d\bm{\lambda}_{n}) \right), \nonumber \\
    & \boxed{ d\bm{y}_{n+1} = (\bm{I}-\tfrac{\Delta{t}}{2}\bm{A})^{-1}(\bm{I}+\tfrac{\Delta{t}}{2}\bm{A})d\bm{y}_{n} + (-\tfrac{\Delta{t}}{2}) (\bm{I}-\tfrac{\Delta{t}}{2}\bm{A})^{-1}\bm{B}d\bm{\lambda}_{n} +
      (-\tfrac{\Delta{t}}{2}) (\bm{I}-\tfrac{\Delta{t}}{2}\bm{A})^{-1}\bm{B}d\bm{\lambda}_{n+1},} \label{m.m.p_eq_1}
\end{align}
where $\bm{A} = H_{\bm{\lambda y}} - H_{\bm{\lambda u}}H_{\bm{uu}}^{-1}H_{\bm{uy}}$, $\bm{B} = H_{\bm{\lambda u}}H_{\bm{uu}}^{-1}H_{\bm{u\lambda}}$.

Applying a similar procedure, we obtain:
\begin{align}
   \boxed{d\bm{\lambda}_{n+1} = (\bm{I}+\tfrac{\Delta{t}}{2}\bm{A}^{T})^{-1}(\bm{I}-\tfrac{\Delta{t}}{2}\bm{A}^{T})d\bm{\lambda}_{n} + (-\tfrac{\Delta{t}}{2}) (\bm{I}+\tfrac{\Delta{t}}{2}\bm{A}^{T})^{-1}\bm{C}d\bm{y}_{n+1} +
      (-\tfrac{\Delta{t}}{2}) (\bm{I}+\tfrac{\Delta{t}}{2}\bm{A}^{T})^{-1}\bm{C}d\bm{y}_{n},} \label{m.m.p_eq_2}
\end{align}
where $\bm{C} = H_{\bm{yy}} - H_{\bm{yu}}H_{\bm{uu}}^{-1}H_{\bm{uy}}$ and
\begin{align*}
    \bm{\lambda}_{n+1} = \bm{\lambda}_{n} - \Delta{t}H_{\bm{y}}(\bm{y_{n+\frac{1}{2}}}, \bm{u}(\bm{y_{n+\frac{1}{2}}},\bm{\lambda_{n+\frac{1}{2}}}),\bm{\lambda_{n+\frac{1}{2}}}).
\end{align*}

By taking the wedge product of \eqref{m.m.p_eq_1} with $d\bm{\lambda}_{n+1}$ we have:
\begin{align}
    d\bm{y}_{n+1} \wedge d\bm{\lambda}_{n+1} =  & (\bm{I}-\tfrac{\Delta{t}}{2}\bm{A})^{-1}(\bm{I}+\tfrac{\Delta{t}}{2}\bm{A})d\bm{y}_{n} \wedge  d\bm{\lambda}_{n+1} + \nonumber \\
    & (-\tfrac{\Delta{t}}{2}) (\bm{I}-\tfrac{\Delta{t}}{2}\bm{A})^{-1}\bm{B}d\bm{\lambda}_{n} \wedge d\bm{\lambda}_{n+1}  = \nonumber \\
    & (\bm{I}-\tfrac{\Delta{t}}{2}\bm{A})^{-1}(\bm{I}+\tfrac{\Delta{t}}{2}\bm{A})d\bm{y}_{n} \wedge (\bm{I}+\tfrac{\Delta{t}}{2}\bm{A}^{T})^{-1}(\bm{I}-\tfrac{\Delta{t}}{2}\bm{A}^{T})d\bm{\lambda}_{n} + \nonumber \\
    & (\bm{I}-\tfrac{\Delta{t}}{2}\bm{A})^{-1}(\bm{I}+\tfrac{\Delta{t}}{2}\bm{A})d\bm{y}_{n} \wedge   (-\tfrac{\Delta{t}}{2}) (\bm{I}+\tfrac{\Delta{t}}{2}\bm{A}^{T})^{-1}\bm{C}d\bm{y}_{n+1} + \nonumber \\
    & (-\tfrac{\Delta{t}}{2}) (\bm{I}-\tfrac{\Delta{t}}{2}\bm{A})^{-1}\bm{B}d\bm{\lambda}_{n} \wedge (-\tfrac{\Delta{t}}{2}) (\bm{I}+\tfrac{\Delta{t}}{2}\bm{A}^{T})^{-1}\bm{C}d\bm{y}_{n+1} + \nonumber \\
    & (-\tfrac{\Delta{t}}{2}) (\bm{I}-\tfrac{\Delta{t}}{2}\bm{A})^{-1}\bm{B}d\bm{\lambda}_{n} \wedge (-\tfrac{\Delta{t}}{2}) (\bm{I}+\tfrac{\Delta{t}}{2}\bm{A}^{T})^{-1}\bm{C}d\bm{y}_{n}. \label{m.m.p_eq_0}
\end{align}
Using the facts $d\bm{a} \wedge d\bm{c} = -d\bm{c} \wedge d\bm{a}$ and $d\bm{a} \wedge \bm{M}d\bm{c} = \bm{M}^{T}d\bm{a} \wedge d\bm{c}$, equation \eqref{m.m.p_eq_0} can be written as:
\begin{align}
    d\bm{y}_{n+1} \wedge d\bm{\lambda}_{n+1} = &  (\bm{I}-\tfrac{\Delta{t}}{2}\bm{A})(\bm{I}+\tfrac{\Delta{t}}{2}\bm{A})^{-1}(\bm{I}-\tfrac{\Delta{t}}{2}\bm{A})^{-1}(\bm{I}+\tfrac{\Delta{t}}{2}\bm{A})  d\bm{y}_{n} \wedge d\bm{\lambda}_{n} + \nonumber \\
    & (\bm{I}-\tfrac{\Delta{t}}{2}\bm{A})^{-1} (\bm{I}+\tfrac{\Delta{t}}{2}\bm{A})^{-1} d\bm{y}_{n} \wedge (-\tfrac{\Delta{t}}{2})(\bm{I}+\tfrac{\Delta{t}}{2}\bm{A}^{T})^{-1}\bm{C} d\bm{y}_{n+1} + \nonumber \\
    & (-\tfrac{\Delta{t}}{2})(\bm{I}-\tfrac{\Delta{t}}{2}\bm{A})^{-1}\bm{B} d\bm{\lambda}_{n} \wedge (-\tfrac{\Delta{t}}{2})(\bm{I}+\tfrac{\Delta{t}}{2}\bm{A}^{T})^{-1}\bm{C} d\bm{y}_{n+1} + \nonumber \\
    & (\tfrac{\Delta{t}}{2})(\bm{I}+\tfrac{\Delta{t}}{2}\bm{A}^{T})^{-1}\bm{C} d\bm{y}_{n} \wedge (-\tfrac{\Delta{t}}{2})(\bm{I}-\tfrac{\Delta{t}}{2}\bm{A})^{-1}\bm{B} d\bm{\lambda}_{n}. \label{m.m.p_eq_3}
\end{align}
Since $(\bm{I} - \tfrac{\Delta{t}}{2}\bm{A})(\bm{I} + \tfrac{\Delta{t}}{2}\bm{A}) = (\bm{I} + \tfrac{\Delta{t}}{2}\bm{A})(\bm{I} - \tfrac{\Delta{t}}{2}\bm{A})$, it follows that $(\bm{I} - \tfrac{\Delta{t}}{2}\bm{A})^{-1}(\bm{I} + \tfrac{\Delta{t}}{2}\bm{A})^{-1} = (\bm{I} + \tfrac{\Delta{t}}{2}\bm{A})^{-1}(\bm{I} - \tfrac{\Delta{t}}{2}\bm{A})^{-1}$, so equation \eqref{m.m.p_eq_3} can be rewritten as:
\begin{align}
     d\bm{y}_{n+1} \wedge d\bm{\lambda}_{n+1} = &  d\bm{y}_{n} \wedge d\bm{\lambda}_{n} + \nonumber \\
    & (\bm{I}-\tfrac{\Delta{t}}{2}\bm{A})^{-1} (\bm{I}+\tfrac{\Delta{t}}{2}\bm{A})^{-1} d\bm{y}_{n} \wedge (-\tfrac{\Delta{t}}{2})(\bm{I}+\tfrac{\Delta{t}}{2}\bm{A}^{T})^{-1}\bm{C} d\bm{y}_{n+1} + \nonumber \\
    & (-\tfrac{\Delta{t}}{2})(\bm{I}-\tfrac{\Delta{t}}{2}\bm{A})^{-1}\bm{B} d\bm{\lambda}_{n} \wedge (-\tfrac{\Delta{t}}{2})(\bm{I}+\tfrac{\Delta{t}}{2}\bm{A}^{T})^{-1}\bm{C} d\bm{y}_{n+1} + \nonumber \\
    & (\tfrac{\Delta{t}}{2})(\bm{I}+\tfrac{\Delta{t}}{2}\bm{A}^{T})^{-1}\bm{C} d\bm{y}_{n} \wedge (-\tfrac{\Delta{t}}{2})(\bm{I}-\tfrac{\Delta{t}}{2}\bm{A})^{-1}\bm{B} d\bm{\lambda}_{n}. \label{m.m.p_eq_4}
\end{align}
\begin{align}
     d\bm{y}_{n+1} \wedge d\bm{\lambda}_{n+1} = & \left(\bm{I}-\tfrac{\Delta{t}^2}{4}\bm{B}^{T}(\bm{I}-\tfrac{\Delta{t}}{2}\bm{A}^{T})^{-1}(\bm{I}+\tfrac{\Delta{t}}{2}\bm{A}^{T})^{-1}\bm{C}\right) d\bm{y}_{n} \wedge d\bm{\lambda}_{n}  + \nonumber \\
     & (\bm{I}-\tfrac{\Delta{t}}{2}\bm{A})^{-1} (\bm{I}+\tfrac{\Delta{t}}{2}\bm{A})^{-1} d\bm{y}_{n} \wedge (-\tfrac{\Delta{t}}{2})(\bm{I}+\tfrac{\Delta{t}}{2}\bm{A}^{T})^{-1}\bm{C} d\bm{y}_{n+1} + \nonumber \\
    & (-\tfrac{\Delta{t}}{2})(\bm{I}-\tfrac{\Delta{t}}{2}\bm{A})^{-1}\bm{B} d\bm{\lambda}_{n} \wedge (-\tfrac{\Delta{t}}{2})(\bm{I}+\tfrac{\Delta{t}}{2}\bm{A}^{T})^{-1}\bm{C} d\bm{y}_{n+1}. \label{m.m.p_eq_5}
\end{align}
By taking the wedge product of $d\bm{y}_{n+1}$ with $(-\tfrac{\Delta{t}}{2})(\bm{I}+\tfrac{\Delta{t}}{2}\bm{A}^{T})^{-1}\bm{C}d\bm{y}_{n+1}$ we have the following:
\begin{align}
   \underbrace{d\bm{y}_{n+1} \wedge  (-\tfrac{\Delta{t}}{2})(\bm{I}+\tfrac{\Delta{t}}{2}\bm{A}^{T})^{-1}\bm{C}d\bm{y}_{n+1}}_{\bm{0}} = & (\bm{I}-\tfrac{\Delta{t}}{2}\bm{A})^{-1}(\bm{I}+\tfrac{\Delta{t}}{2}\bm{A})d\bm{y}_{n} \wedge (-\tfrac{\Delta{t}}{2})(\bm{I}+\tfrac{\Delta{t}}{2}\bm{A}^{T})^{-1}\bm{C}d\bm{y}_{n+1} +\nonumber \\
   & (-\tfrac{\Delta{t}}{2})(\bm{I}-\tfrac{\Delta{t}}{2}\bm{A})^{-1}\bm{B}d\bm{\lambda}_{n} \wedge (-\tfrac{\Delta{t}}{2})(\bm{I}+\tfrac{\Delta{t}}{2}\bm{A}^{T})^{-1}\bm{C}d\bm{y}_{n+1}  + \nonumber \\
   & (-\tfrac{\Delta{t}}{2})(\bm{I}-\tfrac{\Delta{t}}{2}\bm{A})^{-1}\bm{B}d\bm{\lambda}_{n+1} \wedge (-\tfrac{\Delta{t}}{2})(\bm{I}+\tfrac{\Delta{t}}{2}\bm{A}^{T})^{-1}\bm{C}d\bm{y}_{n+1}. \label{m.m.p_eq_6}
\end{align}
Rewriting equation \eqref{m.m.p_eq_6} by isolating $(-\tfrac{\Delta{t}}{2})(\bm{I}-\tfrac{\Delta{t}}{2}\bm{A})^{-1}\bm{B}d\bm{\lambda}_{n+1} \wedge (-\tfrac{\Delta{t}}{2})(\bm{I}+\tfrac{\Delta{t}}{2}\bm{A}^{T})^{-1}\bm{C}d\bm{y}_{n+1}$, we obtain the following equivalent form:
\begin{multline}\label{m.m.p_eq_7}
 (-\tfrac{\Delta{t}}{2})(\bm{I}+\tfrac{\Delta{t}}{2}\bm{A}^{T})^{-1}\bm{C}d\bm{y}_{n+1} \wedge (-\tfrac{\Delta{t}}{2})(\bm{I}-\tfrac{\Delta{t}}{2}\bm{A})^{-1}\bm{B}d\bm{\lambda}_{n+1} = \\
 (\bm{I}-\tfrac{\Delta{t}}{2}\bm{A})^{-1}(\bm{I}+\tfrac{\Delta{t}}{2}\bm{A})d\bm{y}_{n} \wedge (-\tfrac{\Delta{t}}{2})(\bm{I}+\tfrac{\Delta{t}}{2}\bm{A}^{T})^{-1}\bm{C}d\bm{y}_{n+1} + \\
 (-\tfrac{\Delta{t}}{2})(\bm{I}-\tfrac{\Delta{t}}{2}\bm{A})^{-1}\bm{B}d\bm{\lambda}_{n} \wedge (-\tfrac{\Delta{t}}{2})(\bm{I}+\tfrac{\Delta{t}}{2}\bm{A}^{T})^{-1}\bm{C}d\bm{y}_{n+1}.
\end{multline}
Upon plugging \eqref{m.m.p_eq_7} into \eqref{m.m.p_eq_5}, we get:
\begin{multline}\label{m.m.p_eq_8}
    d\bm{y}_{n+1} \wedge d\bm{\lambda}_{n+1} = \left(\bm{I}-\tfrac{\Delta{t}^2}{4}\bm{B}^{T}(\bm{I}-\tfrac{\Delta{t}}{2}\bm{A}^{T})^{-1}(\bm{I}+\tfrac{\Delta{t}}{2}\bm{A}^{T})^{-1}\bm{C}\right) d\bm{y}_{n} \wedge d\bm{\lambda}_{n} \\
    +  \tfrac{\Delta{t}^2}{4}\bm{B}^{T}(\bm{I}-\tfrac{\Delta{t}}{2}\bm{A}^{T})^{-1}(\bm{I}+\tfrac{\Delta{t}}{2}\bm{A}^{T})^{-1}\bm{C} d\bm{y}_{n+1} \wedge d\bm{\lambda}_{n+1}.
\end{multline}
By solving \eqref{m.m.p_eq_8} for $d\bm{y}_{n+1} \wedge d\bm{\lambda}_{n+1}$, we have:
\begin{multline}\label{m.m.p_eq_9}
    \left(\bm{I}-\tfrac{\Delta{t}^2}{4}\bm{B}^{T}(\bm{I}-\tfrac{\Delta{t}}{2}\bm{A}^{T})^{-1}(\bm{I}+\tfrac{\Delta{t}}{2}\bm{A}^{T})^{-1}\bm{C}\right) d\bm{y}_{n+1} \wedge d\bm{\lambda}_{n+1}  = \\
    \left(\bm{I}-\tfrac{\Delta{t}^2}{4}\bm{B}^{T}(\bm{I}-\tfrac{\Delta{t}}{2}\bm{A}^{T})^{-1}(\bm{I}+\tfrac{\Delta{t}}{2}\bm{A}^{T})^{-1}\bm{C}\right) d\bm{y}_{n} \wedge d\bm{\lambda}_{n}.
\end{multline}

\subsection{Numerical results (nonlinear case)}
\label{subsec:numerical-results-nonlinear}
We consider the following nonlinear optimal control problems and apply $\SEoo$, $\SEto$, and the implicit midpoint method. As analytically proven, the condition $H_{yu} = \bm{0}$ in Problem~\eqref{eq_65} ensures the symplecticity of $\SEoo$ and $\SEto$, thus guaranteeing the preservation of the phase-space area. This property is numerically validated in Table~\ref{tab:area-nonlinear}, which compares the absolute relative errors in the preserved area at the initial and final times. In contrast, problem~\eqref{eq_64} does not satisfy the condition $H_{yu} = \bm{0}$, and as shown in Table~\ref{tab:area-nonlinear}, the $\SEoo$ and $\SEto$ methods lose their symplectic property in this case, leading to nonpreservation of the phase-space area.

These simulations employ three distinct initial points to construct a triangular shape at the initial time; see Figure~\ref{fig:triangle}. To compute the area of this triangle, we split it into four smaller non-overlapping triangles, calculate their individual areas, and add them to determine the total area.

\begin{minipage}{0.45\textwidth}
\begin{align}
     \min & \int_{0}^2\, \tfrac{1}{2}v^{2} + \tfrac{\alpha}{2}u^2\, dt \nonumber \\
     & \dot{v} = g + \tfrac{u}{m} - \tfrac{k}{m}v^2 \label{eq_64} \\
     & v(0) = v_0 \nonumber
\end{align}
\end{minipage}
\hfill
\begin{minipage}{0.45\textwidth}
\begin{align}
\min & \int_{0}^2 \, \tfrac{1}{2}v^{2}+ \tfrac{\alpha}{2}u^2\, dt \nonumber \\
& \dot{v} = g + \tfrac{u}{m} - \tfrac{k}{m}uv \label{eq_65} \\
& v(0) = v_0 \nonumber
\end{align}
\end{minipage}

Necessary optimality conditions for \eqref{eq_64} and \eqref{eq_65}, respectively,

\begin{minipage}{0.45\textwidth}
\begin{align}
     & H(v,u,\lambda) = \tfrac{1}{2} v^2 + \tfrac{\alpha}{2} u^2 + \lambda (g + \tfrac{1}{m} u - \tfrac{k}{m}v^2) \nonumber \\
     & \begin{cases}
         \dot{v} = H_{\lambda}(v,u,\lambda) =  g + \tfrac{u}{m} - \tfrac{k}{m}v^2  \\
         \dot{\lambda} =  - H_{v}(v,u,\lambda) = -v + \tfrac{2k}{m} \lambda v  \\
         0 =  H_{u}(v,u,\lambda) = \alpha u + \tfrac{1}{m}\lambda
     \end{cases} \label{eq_66} \\
     & H_{vu} = 0. \nonumber
\end{align}
\end{minipage}
\hfill
\begin{minipage}{0.45\textwidth}
\begin{align}
 & H(v,u,\lambda) = \tfrac{1}{2} v^2 + \tfrac{\alpha}{2} u^2 + \lambda (g + \tfrac{1}{m} u - \tfrac{k}{m} uv) \nonumber \\
 & \begin{cases}
     \dot{v} = H_{\lambda}(v,u,\lambda) =  g + \tfrac{u}{m} - \tfrac{k}{m}uv  \\
     \dot{\lambda} =  - H_{v}(v,u,\lambda) = -v + \tfrac{k}{m} \lambda u  \\
     0 =  H_{u}(v,u,\lambda) = \alpha u + \tfrac{1}{m}\lambda - \tfrac{k}{m}\lambda v
 \end{cases} \label{eq_67} \\
 & H_{vu} \neq 0. \nonumber
\end{align}
\end{minipage}

\begin{table}[ht!]
\centering
\begin{tabular}{|c|c|c|c||c|c|}
\hline
 $\alpha =0.5$, $N.T = 1553$ & $\Delta{t}$ & $|A_0 - A_{T}|$, $H_{vu}=0$ & $\frac{|A_0 - A_{T}|}{A_0}$, $H_{vu} = 0 $ & $|A_0-A_T|$, $H_{vu} \ne \bm{0}$ & $\frac{|A_0 - A_{T}|}{A_0}$, $H_{vu} \neq 0 $ \\ \hline
\multirow{2}{*}{$\SEoo$} & $0.02$ & $3.0978\times 10^{-5}$ & $2.478\times 10^{-2}$ & $1.1293\times 10^{-3}$ & $0.9032$ \\
                         & $0.01$ & $2.2191\times 10^{-5}$ & $1.7753\times 10^{-2}$ & $1.1883\times 10^{-3}$ & $0.9506$ \\ \hline
\multirow{2}{*}{$\SEto$} & $0.02$ & $1.1769\times 10^{-5}$ & $9.415\times 10^{-3}$ & $1.2610\times 10^{-4}$ & $0.1008$ \\
                         & $0.01$ & $1.3989\times 10^{-5}$ & $1.1191\times 10^{-2}$ & $1.5809\times 10^{-4}$ & $0.1264$ \\ \hline
\multirow{2}{*}{MidPoint} & $0.02$ & $1.7307\times 10^{-5}$ & $1.3845\times 10^{-2}$ & $2.0385\times 10^{-4}$ & $0.1630$ \\
                          & $0.01$ & $1.7222\times 10^{-5}$ & $1.3777\times 10^{-2}$ & $1.9555\times 10^{-4}$ & $0.1564$ \\ \hline
\end{tabular}
\caption{Application of $\SEoo$, $\SEto$ and implicit midpoint methods to systems \eqref{eq_66} and \eqref{eq_67} on the interval $[0,2]$ with three initial points $\mathbf{z}_1 = (0,1)^{T}$, $\mathbf{z}_2 =(0.05,1)^{T}$ and $\mathbf{z}_3 = (0,1.05)^{T}$, where $k =1$, $m =20$, $T = 2$. $A_0 = 0.0013$ and $A_T$ denote the areas of the triangle at time $t = 0$ and $t = T = 2$, respectively.}
\label{tab:area-nonlinear}
\end{table}

\begin{table}[ht!]
\centering
\begin{tabular}{|c|c|c|}
   \hline
Method, $\alpha = 0.5$ & $\Delta t$ & $\|A_0^{T} J A_0 - J\|$, $H_{uv}\ne \bm{0}$ \\
   \hline
   \multirow{3}{6em}{$\SEoo$}
   & 0.5  & 0.0174 \\
   & 0.2  & 0.0028 \\
   & 0.02 & $2.7721\times 10^{-5}$ \\
   \hline
   \multirow{3}{6em}{$\SEto$}
   & 0.5  & 0.0166 \\
   & 0.2  & 0.0028 \\
   & 0.02 & $2.7712\times 10^{-5}$ \\
   \hline
   \multirow{3}{6em}{MidPoint}
   & 0.5  & 0.0079 \\
   & 0.2  & $2.7735\times 10^{-4}$ \\
   & 0.02 & $1.5254\times 10^{-7}$ \\
   \hline
   \multirow{3}{6em}{Modified $\SEoo$}
   & 0.5  & $2.5931\times 10^{-8}$ \\
   & 0.2  & $3.6732\times 10^{-9}$ \\
   & 0.02 & $4.9304\times 10^{-10}$ \\
   \hline
   \multirow{3}{6em}{Modified $\SEto$}
   & 0.5  & $1.0876\times 10^{-8}$ \\
   & 0.2  & $3.6732\times 10^{-9}$ \\
   & 0.02 & $4.9304\times 10^{-10}$ \\
   \hline
   \multirow{3}{6em}{Modified MidPoint}
   & 0.5  & $2.4509\times 10^{-8}$ \\
   & 0.2  & $5.0484\times 10^{-9}$ \\
   & 0.02 & $5.1423\times 10^{-10}$ \\
   \hline
\end{tabular}
\caption{Numerical evaluation of symplecticity for $\SEoo$, $\SEto$, $M.P$, $M.\SEoo$, $M.\SEto$, and $M.M.P$ methods (applied to system \eqref{eq_67}) with different time step sizes $(\Delta t)$. $k=1$, $m=20$. $A_0 = \frac{\partial z_{n+1}}{\partial z_n}(v_n,\lambda_n).$}
\label{tab:symp-mod}
\end{table}

\begin{figure}[bht]
    \centering
    \includegraphics[width=0.5\textwidth]{./Images/triangle.pdf}
    \caption{Splitting the triangle with vertices A, B and C into 1553 smaller non-overlapping triangles.}
    \label{fig:triangle}
\end{figure}
\newpage

\section{Nonlinear optimal control problems (Inverted pendulum problem)}
Consider the optimal control problem of the inverted pendulum:
\begin{align*}
    \min_{u}\; &  \tfrac{1}{2}\, \int_{0}^{T} \|\bm{x}-\bm{x}_d\|^{2}\; dt + \tfrac{\alpha}{2} \, \int_{0}^{T} u^{2} dt \\
    & m\ddot{\bm{x}} + \frac{\partial U}{\partial \bm{x}} = \bm{0}, \\
    & \bm{x}(0) = \bm{x}_{0}, \quad \dot{\bm{x}}(0) = \bm{v}_0,
\end{align*}
where
\begin{align*}
    & U = \tfrac{1}{2}k(L-l_0)^2 - \bm{a}^{T}\bm{x},\quad L = \|\bm{x}-(u,0)^{T}\|, \\
    & \frac{\partial U}{\partial\bm{x}} =  k\!\left(\tfrac{L-l_{0}}{L}\right)\!(\bm{x}- (u,0)^{T})-\bm{a}.
\end{align*}
By introducing an auxiliary variable $\bm{v}$ and setting it to $\dot{\bm{x}}$, the second-order differential equation is converted into a first-order system of differential equations, so the control problem is formulated as follows:
\begin{equation}
\label{eq:cont_tim_p}
\begin{aligned}
\min_{\bm{x},\, \bm{v},\, u} & J(\bm{x}, u)
\\
& \dot{\bm{x}} - \bm{v} = \bm{0}, \\
& \dot{\bm{v}} + \frac{1}{m} \frac{\partial U}{\partial\bm{x}} =  \bm{0},
\\
& \bm{x}(0) =\bm{x}_0, \; \bm{v}(0) = \bm{v}_0.
\end{aligned}
\end{equation}

To derive the necessary conditions for optimality, we introduce the Hamiltonian function $H$:
\[
H(\bm{v},\bm{x},u,\bm{\lambda},\bm{\mu}) =
f(\bm{x},u) +  \bm{\lambda}^{T} \!\left(\tfrac{-1}{m}\right)\!\frac{\partial U}{\partial\bm{x}} + \bm{\mu}^{T} \bm{v},
\]
where $f(\bm{x},u) = r(\bm{x}) + q(u)$. By applying Pontryagin's maximum principle to the Hamiltonian function, the necessary conditions for optimality are given as follows:
\begin{align*}
 & \text{Adjoint equations: }
\begin{cases}
\dot{\bm{\lambda}} = -\frac{\partial H}{\partial \bm{v}}, & \bm{\lambda}(T) = \bm{0} \\
\dot{\bm{\mu}} = -\frac{\partial H}{\partial \bm{x}}, & \bm{\mu}(T)= \bm{0}
\end{cases} \\
& \text{Control equation: }
\begin{cases}
\frac{\partial H}{\partial u} = 0
\end{cases} \\
& \text{State equations: }
\begin{cases}
\dot{\bm{v}} = \frac{\partial H}{\partial \bm{\lambda}}, & \bm{v}(0) = \bm{v}_{0} \\
\dot{\bm{x}} = \frac{\partial H}{\partial \bm{\mu}}, & \bm{x}(0) = \bm{x}_{0}
\end{cases}
\end{align*}
This system of equations forms a two-point boundary value problem (TPBVP), with boundary conditions specified at both the initial and terminal times.
\begin{equation}
\label{eq:sys_analytic_sol}
\begin{aligned}
& \begin{cases}
\dot{\bm{\lambda}} + \bm{\mu} = \bm{0}, &  \bm{\lambda}(T) = \bm{0},
\\
\dot{\bm{\mu}} + (\bm{x} - \bm{x}_d)- \left(\tfrac{kl_0}{mL^3}(\bm{x} -(u,0)^{T})(\bm{x}-(u,0)^{T})^{T} + \tfrac{k}{m}\!\left(\tfrac{L-l_0}{L}\right)\bm{I}\right)\bm{\lambda} = \bm{0}, & \bm{\mu}(T) = \bm{0},
\end{cases}
\\
& \alpha u - \bm{\lambda}^{T}\!\left( \left[\tfrac{-kl_0}{mL^{3}}\bm{e}_1^{T}(\bm{x}-(u,0)^{T})\right]\!(\bm{x} - (u,0)^{T}) + \tfrac{-k}{m} \!\left(\tfrac{L-l_0}{L}\right)\bm{e}_1 \right) = 0,
\\
& \begin{cases}
\dot{\bm{v}} + \tfrac{k}{m}\!\left(\tfrac{L-l_0}{L}\right)\!(\bm{x}-(u,0)^{T})-\tfrac{1}{m}\bm{a} = \bm{0}, & \bm{v}(0) = \bm{v}_0,
\\
\dot{\bm{x}} -\bm{v} = \bm{0}, & \bm{x}(0) = \bm{x}_0
\end{cases}
\end{aligned}
\end{equation}
To solve \eqref{eq:sys_analytic_sol}, six numerical schemes are employed as follows:
\begin{align*}
\text{Set } \bm{y}\coloneqq (\bm{v},\bm{x})^{T}, \quad \bm{p} := (\bm{\lambda},\bm{\mu})^{T}.
\end{align*}

\begin{equation}
\begin{array}{c@{\hspace{2em}}c}
SE1: \begin{cases}
            \frac{\Delta\bm{y}}{\Delta t} = \frac{\partial H}{\partial \bm{p}}(\bm{y}_{{\color{red}n+1}}, u_n, \bm{p}_n) \\[4pt]
            \frac{\Delta\bm{p}}{\Delta t} = -\frac{\partial H}{\partial \bm{y}}(\bm{y}_{\color{red}{n+1}}, u_n, \bm{p}_n) \\[4pt]
            H_{u} (\bm{y}_n,u_n,\bm{p}_n)= 0
        \end{cases}  &  M.SE1: \begin{cases}
            \frac{\Delta\bm{y}}{\Delta t} = \frac{\partial H}{\partial \bm{p}}(\bm{y}_{\color{red}n+1}, u_n, \bm{p}_n) \\[4pt]
            \frac{\Delta\bm{p}}{\Delta t} = -\frac{\partial H}{\partial \bm{y}}(\bm{y}_{\color{red}n+1}, u_n, \bm{p}_n) \\[4pt]
            H_{u} (\bm{y}_{\color{red}n+1}, u_{n}, \bm{p}_{n})= 0
        \end{cases} \\[28pt]
SE2: \begin{cases}
            \frac{\Delta\bm{y}}{\Delta t} = \frac{\partial H}{\partial \bm{p}}(\bm{y}_{\color{red}{n}}, u_{n+1}, \bm{p}_{n+1}) \\[4pt]
            \frac{\Delta\bm{p}}{\Delta t} = -\frac{\partial H}{\partial \bm{y}}(\bm{y}_{{\color{red}{n}}}, u_{n+1}, \bm{p}_{n+1}) \\[4pt]
            H_{u}(\bm{y}_{n},u_{n},\bm{p}_{n}) = 0
        \end{cases} &   M.SE2:
        \begin{cases}
            \frac{\Delta\bm{y}}{\Delta t} = \frac{\partial H}{\partial \bm{p}}(\bm{y}_{\color{red}n}, u_{n}, \bm{p}_{n+1}) \\[4pt]
             \frac{\Delta\bm{p}}{\Delta t} = -\frac{\partial H}{\partial \bm{y}}(\bm{y}_{\color{red}n}, u_{n}, \bm{p}_{n+1}) \\[4pt]
             H_{u} (\bm{y}_{\color{red}n}, u_{n}, \bm{p}_{n+1})= 0
        \end{cases} \\[28pt]
        M.P:
        \begin{cases}
            \frac{\Delta\bm{y}}{\Delta t} = \frac{\partial H}{\partial \bm{p}}(\bm{y}_{n+0.5}, u_{n+0.5}, \bm{p}_{n+0.5}) \\[4pt]
             \frac{\Delta\bm{p}}{\Delta t} = -\frac{\partial H}{\partial \bm{y}}(\bm{y}_{n+0.5}, u_{n+0.5}, \bm{p}_{n+0.5}) \\[4pt]
             H_{u} (\bm{y}_n, u_{n}, \bm{p}_{n})= 0
        \end{cases} &  M.M.P:
        \begin{cases}
            \frac{\Delta\bm{y}}{\Delta t} = \frac{\partial H}{\partial \bm{p}}(\bm{y}_{n+0.5}, u_{n}, \bm{p}_{n+0.5}) \\[4pt]
             \frac{\Delta\bm{p}}{\Delta t} = -\frac{\partial H}{\partial \bm{y}}(\bm{y}_{n+0.5}, u_{n}, \bm{p}_{n+0.5}) \\[4pt]
             H_{u} (\bm{y}_{n+0.5}, u_{n}, \bm{p}_{n+0.5})= 0
        \end{cases} \\
\end{array}\label{eq:equation}
\end{equation}
Table~\ref{tab:alpha_methods} presents the numerical results obtained by applying the six numerical schemes for different values of $\alpha$.

\begin{table}[h!]
      \centering
      \begin{tabular}{| l | c | c | c | c | c | c | c | c |}
        \toprule
        \multirow{2}{*}{Method} & \multicolumn{2}{c|}{$\alpha = 10.0$} & \multicolumn{2}{c|}{$\alpha = 5.0$} & \multicolumn{2}{c|}{$\alpha = 1.0$} & \multicolumn{2}{c|}{$\alpha = 0.5$} \\
        \cmidrule{2-9}
        & $\Delta t = 0.1 $& $\Delta t = 0.025 $& $\Delta t = 0.1 $& $\Delta t = 0.025 $& $\Delta t = 0.1 $& $\Delta t = 0.025 $& $\Delta t = 0.1 $& $\Delta t = 0.025 $\\
        \midrule
        SE1 & 6 & 6 & 6 & 6 & 8 & 14 & - & - \\
        M. SE1 & 6 & 6 & 6 & 6 & 7 & 15 & 55 & 104 \\
        \midrule
        SE2 & 6 & 6 & 6 & 6 & 77 & 38 & 84 & 113 \\
        M. SE2 & 6 & 6 & 6 & 6 & 26 & 28 & 29 & 41 \\
        \midrule
        MidPoint & 6 & 6 & 6 & 6 & 25 & 24 & 67 & 51 \\
        M. MidPoint & 6 & 6 & 6 & 6 & 19 & 19 & 29 & 61 \\
        \bottomrule
      \end{tabular}
      \caption{Number of iterations (N.I.) for different methods and values of \(\alpha\), with $\Delta t \in \{0.1, 0.025\}$, $k=1$, $m=1$, $\bm{x}_{0} = (1,1)^{T}$, $\bm{x}_{d} = (0,2)^{T}$, $\bm{v}_{0} =(0,0)^{T}$, $l_{0} = \sqrt{2}$, $T \in [0,10]$.}
      \label{tab:alpha_methods}
    \end{table}

\begin{table}[h!]
      \centering
      \caption{Numerical evaluation of symplecticity for SE1, SE2, M.P, M.SE1, M.SE2, and M.M.P methods with different time step sizes $\Delta t$. $\|\bm{A}_0^{T}\bm{J}\bm{A}_0-\bm{J}\|$, where $\bm{A}_{0} = \frac{\partial(\bm{y}_{n+1},\bm{p}_{n+1})}{\partial(\bm{y}_{n},\bm{p}_{n})}$.}
      \label{tab:symp-pendulum}
      \vspace{4pt}
      \resizebox{\textwidth}{!}{%
      \begin{tabular}{| l | c | c | c | c | c | c | c | c | c |}
        \toprule
        \multirow{2}{*}{Method}
          & \multicolumn{3}{c|}{$\alpha = 10.0$}
          & \multicolumn{3}{c|}{$\alpha = 5.0$}
          & \multicolumn{3}{c|}{$\alpha = 1.0$} \\
        \cmidrule{2-10}
          & $\Delta t=1$ & $\Delta t=0.1$ & $\Delta t=0.025$
          & $\Delta t=1$ & $\Delta t=0.1$ & $\Delta t=0.025$
          & $\Delta t=1$ & $\Delta t=0.1$ & $\Delta t=0.025$ \\
        \midrule
        SE1 & 0.0614853 & 0.0030918 & $1.4210\times 10^{-4}$
            & 0.1148578 & 0.0031728 & $1.0397\times 10^{-4}$
            & 0.3654304 & 0.0045715 & $1.7756\times 10^{-4}$ \\
        M.\ SE1 & $3.7994\times 10^{-10}$ & $4.7898\times 10^{-10}$ & $6.8444\times 10^{-10}$
                & $3.8148\times 10^{-10}$ & $3.9087\times 10^{-10}$ & $6.7330\times 10^{-10}$
                & $3.9091\times 10^{-10}$ & $4.9181\times 10^{-10}$ & $4.0767\times 10^{-10}$ \\
        \midrule
        SE2 & 81.5165 & $5.5104\times 10^{-4}$ & $8.2952\times 10^{-5}$
            & 29.9389 & 0.0021965 & $3.2923\times 10^{-5}$
            & 25.2454 & 0.1528280 & $5.3776\times 10^{-5}$ \\
        M.\ SE2 & $4.1198\times 10^{-8}$ & $1.2966\times 10^{-9}$ & $5.8667\times 10^{-10}$
                & $4.3593\times 10^{-8}$ & $6.6569\times 10^{-10}$ & $4.4809\times 10^{-10}$
                & $1.9479\times 10^{-8}$ & $5.4068\times 10^{-10}$ & $7.0251\times 10^{-10}$ \\
        \midrule
        MidPoint & 0.1192 & 0.0020301 & $1.3219\times 10^{-4}$
                & 0.1403 & 0.0026914 & $1.7456\times 10^{-4}$
                & 0.4556 & 0.0121969 & $7.6003\times 10^{-4}$ \\
        M.\ MidPoint & $1.3738\times 10^{-9}$ & $4.6400\times 10^{-10}$ & $7.4220\times 10^{-10}$
                    & $1.0867\times 10^{-9}$ & $5.4220\times 10^{-10}$ & $5.9989\times 10^{-10}$
                    & $6.6610\times 10^{-10}$ & $6.0072\times 10^{-10}$ & $7.4891\times 10^{-10}$ \\
        \bottomrule
      \end{tabular}
      }
\end{table}

\pagebreak
\section{Algorithm Comparison}

In this section, we compare the baseline algorithms (SE1, SE2, MidPoint) with their modified counterparts.
The comparison is illustrated using three types of plots:

\begin{itemize}
  \item \textbf{Iteration plots:} visualize convergence behavior over a grid of initial states \(\mathbf{x_0}\).
  For each tested \(\mathbf{x_0}\), the color encodes the number of iterations required to converge (up to a fixed maximum).
  If the method does not converge within the maximum number of iterations (30 in our experiments), the corresponding point is colored black.
  Additionally, an arrow is drawn from each \(\mathbf{x_0}\) to the final state \(\mathbf{x_N}\), indicating the achieved terminal state.

  \item \textbf{Control plots:} show the computed control sequence \(u(t_i)\) (or \(u_i\)) over the discrete time index $t_i$, where \(i \in \{0,\dots,N\}\).

  \item \textbf{Trajectory plots:} show the path of the system in state space as it evolves from \(\mathbf{x_0}\) to \(\mathbf{x_N}\).
  Unlike the iteration plot arrows, these plots display the full trajectory, not only the start and end points.
\end{itemize}

\noindent Before presenting the full set of results, we first explain these plots using a representative example.

\begin{figure}[H]
  \centering
  \includegraphics[width=\textwidth]{Images/IterationPlot_Modified SE1_k=0.25_h=10.0_N=200_m=1.0}
  \caption{Iteration plot of Modified SE1 for \(k=0.25\), \(h=10.0\), \(N=200\), and \(m=1.0\).}
  \label{fig:example_iter_plot}
\end{figure}

\noindent Figure~\ref{fig:example_iter_plot} shows an iteration plot for the \emph{Modified SE1} method with
\(k=0.25\), \(h=10.0\), \(N=200\), \(m=1.0\), and desired terminal state
\(\mathbf{x_d}=\vecx{0.00}{10.00}\).
Each dot corresponds to a different initial condition \(\mathbf{x_0}\) in the sampled grid.
The color indicates the number of iterations needed to converge; points that do not converge within the allowed maximum are colored black.
Finally, each arrow maps \(\mathbf{x_0}\) to the achieved terminal state \(\mathbf{x_N}\).
Note that these arrows do \emph{not} represent the full trajectory through the state space---only the start and end points.
The full state evolution is shown in the trajectory plot.

\pagebreak

\begin{figure}[H]
  \centering
  \includegraphics[width=\textwidth]{Images/ControlPlot_Modified SE1_k=0.25_h=10.0_N=200_m=1.0_x0=-1.11, 2.12}
  \caption{Control plot of Modified SE1 for \(k=0.25\), \(h=10.0\), \(N=200\), \(m=1.0\), and \(\mathbf{x_0}=\vecx{-1.11}{2.12}\).}
  \label{fig:example_control_plot}
\end{figure}

\noindent Figure~\ref{fig:example_control_plot} shows the computed control sequence for Modified SE1 with
\(k=0.25\), \(h=10.0\), \(N=200\), \(m=1.0\), desired state \(\mathbf{x_d}=\vecx{0.00}{10.00}\),
and initial condition \(\mathbf{x_0}=\vecx{-1.11}{2.12}\).
For modified methods, the horizontal axis corresponds to discrete time samples \(t_i\), where the exact interpretation depends on the modification.
In practice, the sampling approaches the uniform grid as \(N \to \infty\), i.e.,
\[
t_i \to t_0 + \frac{T-t_0}{N}\, i.
\]
For the unmodified methods, the grid is explicitly uniform:
\[
t_i = t_0 + \frac{T-t_0}{N}\, i,
\qquad i \in \{0,\dots,N\}.
\]

\begin{figure}[H]
  \centering
  \includegraphics[width=0.8\textwidth]{Images/ProjectionPlot_Modified SE1_k=0.25_h=10.0_N=200_m=1.0_x0=-1.11, 2.12}
  \caption{Trajectory plot of Modified SE1 for \(k=0.25\), \(h=10.0\), \(N=200\), \(m=1.0\), and \(\mathbf{x_0}=\vecx{-1.11}{2.12}\).}
  \label{fig:example_projection_plot}
\end{figure}

\noindent Figure~\ref{fig:example_projection_plot} shows the resulting trajectory in state space for the same configuration as in Figure~\ref{fig:example_control_plot}.
In particular, the plot visualizes the evolution of the state from \(\mathbf{x_0}=\vecx{-1.11}{2.12}\) to the final state \(\mathbf{x_N}\),
which corresponds to the endpoint indicated in the iteration plot (Figure~\ref{fig:example_iter_plot}).

\medskip
\noindent Finally, note that because the number of generated examples is large, we only include representative cases in which the difference between baseline and modified trajectories is pronounced (i.e., where the trajectory discrepancy is high).
Moreover, when the desired height component of \(\mathbf{x_d}\) is small, the converged solutions often correspond to trajectories in which the pendulum falls.
Empirically, such cases can exhibit large differences in the projection plots even when the terminal states \(\mathbf{x_N}\) are close.

\FloatBarrier
\subsection{SE1}
\algoAnalysis[
  method = SE1,
  k      = 1.0,
  h      = 8.0,
  N      = 200,
  m      = 1.0,
  xA     = {2.00, 2.93},
  xB     = {-1.56, 2.93}
]

\algoAnalysis[
  method = SE1,
  k      = 0.5,
  h      = 8.0,
  N      = 100,
  m      = 1.0,
  xA     = {-2.00, 1.72},
]

\algoAnalysis[
    method = SE1,
    k      = 0.25,
    h      = 10.0,
    N      = 100,
    m      = 1.0,
    xA     = {1.56, 1.72},
]

\subsection{SE2}
\algoAnalysis[
    method = SE2,
    k      = 1.0,
    h      = 6.0,
    N      = 100,
    m      = 1.0,
    xA     = {1.11, 2.52},
]

\algoAnalysis[
    method = SE2,
    k      = 0.5,
    h      = 10.0,
    N      = 100,
    m      = 1.0,
    xA     = {-2.00, 3.74},
    xB     = {2.00, 3.74},
]

\algoAnalysis[
    method = SE2,
    k      = 0.25,
    h      = 10.0,
    N      = 100,
    m      = 1.0,
    xA     = {-1.56, 2.12},
]

\subsection{MidPoint}
\algoAnalysis[
    method = MidPoint,
    k      = 1.0,
    h      = 10.0,
    N      = 200,
    m      = 1.0,
    xA     = {1.11, 3.74},
    xB     = {2.00, 3.74},
    xC     = {-2.00, 3.74},
]

\algoAnalysis[
    method = MidPoint,
    k      = 0.5,
    h      = 10.0,
    N      = 100,
    m      = 1.0,
    xA     = {-1.11, 3.33},
]

\algoAnalysis[
    method = MidPoint,
    k      = 0.25,
    h      = 10.0,
    N      = 100,
    m      = 1.0,
    xA     = {0.22, 2.12},
]

    \section{Forward-Backward Euler and RK Sweep Schemes}
\label{sec:fb-as-guess}

Newton's method, when supplied with a sufficiently good initial guess,
converges to the solution of the discrete optimality conditions in only a
handful of iterations. In our setting this is a substantial advantage: each
Newton step operates monolithically on the full discretised
two-point boundary value problem and is comparatively expensive, but the
per-iteration progress is large. The bottleneck is therefore not the cost
of an individual step but the requirement that the iterate already lie
inside the basin of attraction of Newton's method. From a random or
otherwise poorly chosen starting point, convergence is unreliable and
frequently fails outright.

For the inverted-pendulum problem with $\alpha = 10$ we were able to
construct an educated guess by hand. The intuition is straightforward: the
regularisation is strong enough that the optimal control keeps the
pendulum from falling, pushing the ball laterally so that it eventually
reaches the $y$-axis, after which the pendulum ascends along a near-linear
path toward the origin. Encoding this picture into an initial guess
produced reliable Newton convergence for $\alpha = 10$ across most of the
sampled initial states.

This heuristic, however, is fragile. As $\alpha$ decreases, the
regularisation weakens and the optimal trajectory deviates substantially
from the linear-ascent picture; the educated guess no longer lies in
Newton's basin and the method stalls or diverges. The same failure mode
appears at fixed $\alpha$ when the initial state $\mathbf{x_0}$ corresponds
to a hard configuration --- a low height or an awkward angle --- even when
$\alpha$ is large.

To address this, we use the forward--backward gradient method to generate
a ``good enough'' initial guess, which Newton's method then refines.
Forward--backward sweeps are cheap on a per-iteration basis: each sweep
propagates the state forward, the adjoint backward, and updates the
control without forming or factorising a Jacobian of the full discretised
system. The trade-off is the opposite of Newton's: many cheap iterations
rather than a few expensive ones. Because the gradient method only needs
to deliver an iterate inside Newton's basin --- not a tightly converged
solution --- it is well-suited to this preconditioning role.

We computed full forward--backward solutions with tolerance
$\varepsilon = 10^{-12}$ using two variants:
\begin{itemize}
    \item the classical forward--backward sweep with explicit Euler
    propagation of the state and adjoint, and
    \item an RK4 forward--backward sweep, which uses fourth-order
    Runge--Kutta integration for both equations.
\end{itemize}
All examples share $\mathbf{x_d} = (0.0,\,10.0)^{T}$, $m = 1.0$, $k = 1.0$,
$l_0 = 9.0$, and time horizon $[0,\,10]$\,s. We swept
$\alpha \in \{10,\,5,\,1\}$ over the initial-state grid
$\mathbf{x_0} \in \{1\}\times\{1,\,2,\,3,\,4\}$, with grid sizes
$N \in \{100,\,200,\,400,\,600,\,800\}$. The study is currently
being extended to $\alpha \in \{0.5,\,0.1\}$.

For each configuration we report three views of the converged solution:
the Hamiltonian $H(t_i)$, the control $u(t_i)$, and the residual cost as a
function of iteration index. Together they expose how the solution refines
as $N$ increases: $H$ becomes flatter, reflecting its first-integral
character along the optimal trajectory; $u$ becomes smoother; and the
residual decays more cleanly.

\paragraph{Number of iterations and convergence cap.}
We recorded both the number of iterations and the wall time of each
method. The number of iterations can be large --- ranging from a few
hundred for the easiest cases (e.g.\ $\alpha = 10$, $\mathbf{x_0} =
(1,4)^T$, $N = 100$: $444$ RK4-FB iterations, $749$ classical-FB
iterations) to several hundred thousand for the hardest --- but each
iteration is cheap, so the overall wall time remains tractable. To
bound runtime we capped the gradient iteration at $500{,}000$. A
recorded number of iterations equal to $500{,}001$ in the tables
therefore signals that the run hit the cap rather than meeting the
tolerance.

A non-trivial fraction of the runs in our grid fail to reach
tolerance, predominantly at small $\alpha$ and large $N$. The non-convergence
is visible both in the number of iterations --- many runs hit the cap of
$500{,}000$ iterations --- and in the diagnostic plots: the Hamiltonian and
control profiles fail to settle to a clean shape, exhibiting noise and
oscillations absent from the converged cases. The mechanism behind these
failures is a collapse of the step-length rule used in the gradient
update, which we describe in detail in
Section~\ref{sec:fb-collapse}, and the per-run outcome of every
$(\alpha, N, \bm{x}_0)$ in our grid is tabulated in
Table~\ref{tab:fb-results}.

\subsection{Discrete forward--backward sweep schemes}
\label{sec:fb-schemes}

The numerical results of Section~\ref{sec:fb-numerical} are produced
by two discrete forward--backward sweeps, identical in outer structure
and differing only in how the state and adjoint ODEs are integrated.
Both schemes target the following optimal-control problems
\begin{equation}
    \min_{\bm{u}}\; J(\bm{u}) = \int_{a}^{b} L(t,\bm{y},\bm{u})\,dt
    \quad \text{s.t.} \quad \dot{\bm{y}} = \bm{f}(t,\bm{y},\bm{u}),
    \quad \bm{y}(a) = \bm{y}_0,
\end{equation}
which specialises to the inverted-pendulum problem
\eqref{eq:sys_analytic_sol} on identification of $\bm{y} \coloneqq
(\bm{v},\bm{x})^{T}$ and $\bm{p} \coloneqq (\bm{\lambda},\bm{\mu})^{T}$.
With Hamiltonian
\begin{equation}
    H(t,\bm{y},\bm{u},\bm{p}) = L(t,\bm{y},\bm{u}) + \bm{p}^{T}\bm{f}(t,\bm{y},\bm{u}),
\end{equation}
the continuous optimality system reads
\begin{equation}
    \dot{\bm{y}} = \bm{f},
    \qquad
    \dot{\bm{p}} = -H_{\bm{y}} = -L_{\bm{y}} - \bm{f}_{\bm{y}}^{T}\bm{p},
    \quad \bm{p}(b) = \bm{0},
    \qquad
    \nabla J(\bm{u})(t) = H_{\bm{u}} = L_{\bm{u}} + \bm{f}_{\bm{u}}^{T}\bm{p}.
    \label{eq:fb-cont-opt}
\end{equation}
The interval $[a,b]$ is divided into $N$ equal subintervals of length
$h = (b-a)/N$, with nodes $t_n = a + n h$, $n = 0,\dots,N$. We write
$\bm{y}_n \approx \bm{y}(t_n)$, $\bm{u}_n \approx \bm{u}(t_n)$,
$\bm{p}_n \approx \bm{p}(t_n)$, and use the superscript $(k)$ to denote
the gradient-descent iteration index.

\subsubsection{Forward--backward sweep with explicit Euler}
\label{sec:fb-euler}

For the Euler variant the state is advanced forward by an explicit Euler
step,
\begin{equation}
    \bm{y}_0^{(k)} = \bm{y}_0, \qquad
    \bm{y}_{n+1}^{(k)} = \bm{y}_n^{(k)} + h\,\bm{f}\!\bigl(t_n,\, \bm{y}_n^{(k)},\, \bm{u}_n^{(k)}\bigr),
    \quad n = 0,\dots,N-1.
    \label{eq:fb-eul-state}
\end{equation}
Because the Euler update at step $n$ depends on $\bm{u}_n$ only, the value
$\bm{u}_N$ does not enter the state sweep; in practice the control is
therefore stored at $t_0,\dots,t_{N-1}$.

For the adjoint, $\dot{\bm{p}} = -H_{\bm{y}}$ is discretised backward in
time. Writing the backward Euler step as
$\bm{p}_{n-1} = \bm{p}_n - h\,\dot{\bm{p}}(t_n) = \bm{p}_n + h\,H_{\bm{y}}(t_n,\bm{y}_n,\bm{u}_n,\bm{p}_n)$
--- the plus sign coming from $\dot{\bm{p}} = -H_{\bm{y}}$ --- the discrete
adjoint sweep is
\begin{equation}
    \bm{p}_N^{(k)} = \bm{0}, \qquad
    \bm{p}_{n-1}^{(k)} = \bm{p}_n^{(k)}
        + h\bigl[L_{\bm{y}}(t_n,\bm{y}_n^{(k)},\bm{u}_n^{(k)})
        + \bm{f}_{\bm{y}}(t_n,\bm{y}_n^{(k)},\bm{u}_n^{(k)})^{T}\bm{p}_n^{(k)}\bigr],
    \quad n = N,\dots,1.
    \label{eq:fb-eul-adjoint}
\end{equation}

The discrete gradient at node $t_n$ is
\begin{equation}
    \bm{g}_n^{(k)} \coloneqq H_{\bm{u}}\bigl(t_n,\bm{y}_n^{(k)},\bm{u}_n^{(k)},\bm{p}_n^{(k)}\bigr)
    = L_{\bm{u}}(t_n,\bm{y}_n^{(k)},\bm{u}_n^{(k)})
        + \bm{f}_{\bm{u}}(t_n,\bm{y}_n^{(k)},\bm{u}_n^{(k)})^{T}\bm{p}_n^{(k)},
    \label{eq:fb-grad}
\end{equation}
and the control is updated by gradient descent,
\begin{equation}
    \bm{u_n}^{(k+1)} = \bm{u_n}^{(k)} - \alpha_k\,\bm{g_n}^{(k)}
    \label{eq:fb-ctrl-update}
\end{equation}
with step length $\alpha_k$ chosen as in Section~\ref{sec:fb-stabbb}. The
iteration is started from an initial guess
$\{\bm{u}_n^{(0)}\}_{n=0}^{N-1}$ and stopped when either
\begin{equation}
    \bigl\|\bm{u}^{(k+1)} - \bm{u}^{(k)}\bigr\| < \varepsilon
    \quad \text{or} \quad
    \bigl\|\bm{g}^{(k)}\bigr\| < \varepsilon,
    \label{eq:fb-stop}
\end{equation}
where $\bm{u}^{(k)}$ and $\bm{g}^{(k)}$ are the stacked control and
gradient vectors at iteration $k$ (over all nodes), and $\|\cdot\|$ is
the Euclidean norm.
The role of the \emph{or} is discussed in Section~\ref{sec:fb-collapse}.

\subsubsection{Forward--backward sweep with RK4}
\label{sec:fb-rk4}

The RK4 variant replaces both Euler integrations by classical fourth-order
Runge--Kutta. Midpoint quantities are required both for the forward sweep
at $t_{n+\tfrac12} = t_n + h/2$ and for the backward sweep at
$t_{n-\tfrac12} = t_n - h/2$; we approximate the control and state at
midpoints by linear interpolation between the surrounding nodes,
\begin{equation}
    \bm{u}_{n+\tfrac12} \approx \tfrac{\bm{u}_n + \bm{u}_{n+1}}{2},
    \qquad
    \bm{u}_{n-\tfrac12} \approx \tfrac{\bm{u}_{n-1} + \bm{u}_n}{2},
    \qquad
    \bm{y}_{n-\tfrac12} \approx \tfrac{\bm{y}_{n-1} + \bm{y}_n}{2}.
    \label{eq:fb-rk4-mid}
\end{equation}
Unlike the Euler variant, the RK4 forward update at step $n$ involves
$\bm{u}_{n+1}$ through $\bm{k}_4$ (see below), so the control is stored at
\emph{every} node $t_0,\dots,t_N$. The backward sweep additionally reuses
the state midpoint $\bm{y}_{n-\tfrac12}$ from the converged forward
trajectory, rather than recomputing intermediate stages.

\paragraph{Forward sweep.} For $n = 0,\dots,N-1$:
\begin{align}
    \bm{k}_1 &= \bm{f}\bigl(t_n,\,\bm{y}_n^{(k)},\,\bm{u}_n^{(k)}\bigr), \\
    \bm{k}_2 &= \bm{f}\bigl(t_n + \tfrac{h}{2},\,\bm{y}_n^{(k)} + \tfrac{h}{2}\bm{k}_1,\,\bm{u}_{n+\tfrac12}^{(k)}\bigr), \\
    \bm{k}_3 &= \bm{f}\bigl(t_n + \tfrac{h}{2},\,\bm{y}_n^{(k)} + \tfrac{h}{2}\bm{k}_2,\,\bm{u}_{n+\tfrac12}^{(k)}\bigr), \\
    \bm{k}_4 &= \bm{f}\bigl(t_n + h,\,\bm{y}_n^{(k)} + h\bm{k}_3,\,\bm{u}_{n+1}^{(k)}\bigr), \\
    \bm{y}_{n+1}^{(k)} &= \bm{y}_n^{(k)} + \tfrac{h}{6}\bigl(\bm{k}_1 + 2\bm{k}_2 + 2\bm{k}_3 + \bm{k}_4\bigr),
    \label{eq:fb-rk4-fwd}
\end{align}
with $\bm{y}_0^{(k)} = \bm{y}_0$.

\paragraph{Backward sweep.} Defining
\begin{equation}
    g(t,\bm{y},\bm{u},\bm{p}) \coloneqq -L_{\bm{y}}(t,\bm{y},\bm{u}) - \bm{f}_{\bm{y}}(t,\bm{y},\bm{u})^{T}\bm{p}
\end{equation}
so that $\dot{\bm{p}} = g$, the RK4 stages for $n = N, N-1,\dots,1$ read
\begin{align}
    \bm{\ell}_1 &= g\bigl(t_n,\,\bm{y}_n^{(k)},\,\bm{u}_n^{(k)},\,\bm{p}_n^{(k)}\bigr), \\
    \bm{\ell}_2 &= g\bigl(t_n - \tfrac{h}{2},\,\bm{y}_{n-\tfrac12}^{(k)},\,\bm{u}_{n-\tfrac12}^{(k)},\,\bm{p}_n^{(k)} - \tfrac{h}{2}\bm{\ell}_1\bigr), \\
    \bm{\ell}_3 &= g\bigl(t_n - \tfrac{h}{2},\,\bm{y}_{n-\tfrac12}^{(k)},\,\bm{u}_{n-\tfrac12}^{(k)},\,\bm{p}_n^{(k)} - \tfrac{h}{2}\bm{\ell}_2\bigr), \\
    \bm{\ell}_4 &= g\bigl(t_{n-1},\,\bm{y}_{n-1}^{(k)},\,\bm{u}_{n-1}^{(k)},\,\bm{p}_n^{(k)} - h\bm{\ell}_3\bigr), \\
    \bm{p}_{n-1}^{(k)} &= \bm{p}_n^{(k)} - \tfrac{h}{6}\bigl(\bm{\ell}_1 + 2\bm{\ell}_2 + 2\bm{\ell}_3 + \bm{\ell}_4\bigr),
    \label{eq:fb-rk4-bwd}
\end{align}
with $\bm{p}_N^{(k)} = \bm{0}$.

The gradient is evaluated at every node $n = 0,\dots,N$ via
\eqref{eq:fb-grad}; the control update \eqref{eq:fb-ctrl-update} is
applied at every node $n = 0,\dots,N$; and the stopping rule
\eqref{eq:fb-stop} is identical to the Euler variant.

\subsection{Step length: stabilised Barzilai--Borwein}
\label{sec:fb-stabbb}

The step length $\alpha_k$ in \eqref{eq:fb-ctrl-update} is critical to
convergence: too small and the iteration crawls, too large and the cost
diverges. Throughout this work $\alpha_k$ is set by the stabilised
Barzilai--Borwein (BB) line search of Bijalwan and Mu\~noz (Computer
Methods in Applied Mechanics and Engineering~420, 2024, Algorithm~2),
which combines the short BB step with a safeguard against overlong
increments.

Stacking the control and gradient over all nodes into the vectors
$\bm{u}^{k} = (\bm{u}_n^{(k)})_n$ and $\bm{d}^{k} = (\bm{g}_n^{(k)})_n$,
let $\Delta\bm{u} = \bm{u}^{k} - \bm{u}^{k-1}$ and
$\Delta\bm{d} = \bm{d}^{k} - \bm{d}^{k-1}$. The short BB step is
\begin{equation}
    \theta_s^{k} = -\,\frac{\Delta\bm{u}^{T}\Delta\bm{d}}{\|\Delta\bm{d}\|^{2}}.
    \label{eq:bb-short}
\end{equation}
For non-convex objectives $\theta_s^{k}$ may be negative; the geometric-mean
fallback
\begin{equation}
    \theta_m^{k} = \frac{\|\Delta\bm{u}\|}{\|\Delta\bm{d}\|}
    \label{eq:bb-mean}
\end{equation}
is then used instead. To prevent overlong increments, an upper bound
$\theta_{\max} > 0$ on the displacement $\|\bm{u}^{k+1}-\bm{u}^{k}\|$ is
imposed, which yields the threshold
\begin{equation}
    \theta_{\mathrm{th}}^{k} = \frac{\theta_{\max}}{\|\bm{d}^{k}\|}.
    \label{eq:bb-thresh}
\end{equation}
The final BB step length used in \eqref{eq:fb-ctrl-update} is
\begin{equation}
    \alpha_k = \theta^{k} =
    \begin{cases}
        \min\bigl(\theta_s^{k},\,\theta_{\mathrm{th}}^{k}\bigr), & \theta_s^{k} > 0, \\
        \min\bigl(\theta_m^{k},\,\theta_{\mathrm{th}}^{k}\bigr), & \theta_s^{k} < 0,
    \end{cases}
    \label{eq:bb-final}
\end{equation}
with $\theta_{\max}$ a fixed positive parameter of the algorithm. The
Stabilised BB rule \eqref{eq:bb-final} requires the previous increments
$\Delta\bm{u}, \Delta\bm{d}$, which are only defined for $k \geq 1$; we
therefore supply $\alpha_0$ as an externally specified input parameter
and apply \eqref{eq:bb-final} from $k = 1$ onward.

\subsection{Premature termination from BB collapse}
\label{sec:fb-collapse}

The stabilised BB rule of Section~\ref{sec:fb-stabbb} bounds $\alpha_k$
from above but \emph{not} from below. Near a stationary point the
increments $\Delta\bm{u}$ and $\Delta\bm{d}$ shrink together at comparable
rates, so the ratios in \eqref{eq:bb-short}--\eqref{eq:bb-mean} approach
$0/0$ and eventually return $\mathrm{NaN}$ in floating-point arithmetic.
Once $\alpha_k = \mathrm{NaN}$, the next control update
\eqref{eq:fb-ctrl-update} destroys the iterate.

The collapse is most severe at small $\alpha$ --- the regularisation
weight in the cost, not to be confused with the BB step $\alpha_k$ ---
where the gradient
$H_{\bm{u}} = \alpha\,\bm{u} + \dots$ responds weakly to changes in
$\bm{u}$, and at large $N$, where per-component magnitudes shrink with
the time step. In those regimes the BB step length collapses toward zero
faster than $\|H_{\bm{u}}\|$ falls below the tolerance $\varepsilon$.

To avoid silent corruption of the iterate we relaxed the convergence
test from the conjunction
\begin{equation*}
    \bigl\|\bm{u}^{(k+1)} - \bm{u}^{(k)}\bigr\| < \varepsilon
    \quad \text{and} \quad
    \bigl\|H_{\bm{u}}^{(k)}\bigr\| < \varepsilon
\end{equation*}
to the disjunction in \eqref{eq:fb-stop}: the iteration stops as soon as
\emph{either} the control update or the gradient becomes small. The
price is that some runs terminate while $\|H_{\bm{u}}\|$ is still well
above $\varepsilon$; these are the rows shaded in
Table~\ref{tab:fb-results}. The final BB step lengths
$\theta^{k}_{\mathrm{last}}$ in those rows are several orders of
magnitude smaller than in the converged rows at the same
$(\alpha, N, \bm{x}_0)$, which is precisely the BB-collapse signature.

\section{Forward--Backward Numerical Results}
\label{sec:fb-numerical}

We present the converged forward--backward solutions across the
$(\Delta t,\,\alpha,\,\mathbf{x_0})$ grid described in
Section~\ref{sec:fb-as-guess}. For each configuration we report two views
of the converged solution: the Hamiltonian
$H_i \coloneqq H(\mathbf{y}_i,u_i,\mathbf{p}_i)$ along the discretised
trajectory and the optimal control $u_i$. Each plot compares the two
forward--backward variants on the same axes: classical Euler
forward--backward (blue) and RK4 forward--backward (orange).

We organise the results by $\alpha$. For each $(\alpha, \mathbf{x_0})$
shown we display the Hamiltonian and the control as $\Delta t$ is
decreased from $0.1$ down to $0.0125$. We restrict the figures to
configurations in which \emph{every} $\Delta t$ in the sweep converged
according to the criterion of Section~\ref{sec:fb-collapse}; the full
convergence outcome of every $(\alpha, N, \mathbf{x_0})$ in our grid is
in Table~\ref{tab:fb-results}. With this restriction, only $\alpha = 10$
shows the complete $\mathbf{x_0} \in \{(1,1),(1,2),(1,3),(1,4)\}^{T}$
sweep; at $\alpha = 5$ only $\mathbf{x_0} = (1,2)^{T}$ converged across
every $\Delta t$; and at $\alpha = 1$ no initial state did. Two trends
are universal across the converged cases: as $\Delta t$ decreases, the
Hamiltonian flattens toward the constant value expected of a first
integral of the optimality system, and the control profile sharpens
around the boundary layers and stabilises in the bulk. The discrepancies
between Euler and RK4 forward--backward shrink with $\Delta t$ as well,
since both variants are recovering the same continuous-time optimum.

\FloatBarrier
\subsection{Convergence and non-convergence}
\label{sec:fb-convergence}

In a non-trivial fraction of configurations the gradient method failed to
reach tolerance. The diagnostic plots make these runs unmistakable: $H_i$
is far from constant and the control profile fails to settle, instead
exhibiting persistent oscillations and discontinuous jumps. The mechanism
behind the failure is the BB step-length collapse analysed in
Section~\ref{sec:fb-collapse}: the iteration is stopped by the
$\|\Delta\bm{u}\| < \varepsilon$ branch of \eqref{eq:fb-stop} while
$\|H_{\bm{u}}\|$ is still well above $\varepsilon$, because the
stabilised BB rule has driven $\theta^{k}$ to floating-point underflow.
Such runs are not used as Newton seeds.

Table~\ref{tab:fb-results} reports, for every
$(\alpha, N, \bm{x}_0)$ in our grid, the number of iterations, wall time,
final $\|H_{\bm{u}}\|$, and final BB step length
$\theta^{k}_{\mathrm{last}}$ for both Euler and RK4 forward--backward.
Runs flagged as failed --- those where $\|H_{\bm{u}}\| > 10^{-5}$ at
termination --- are shaded; their $\theta^{k}_{\mathrm{last}}$ entries
are several orders of magnitude smaller than the converged rows, which
is the BB-collapse signature described in
Section~\ref{sec:fb-collapse}. A dash ``\textbf{---}'' in the number-of-iterations
column means the run hit the $500{,}000$-iteration cap before any
stopping criterion fired. The figures below restrict to
configurations in which every $\Delta t$ converged.

\input{tables/fb_results_table}

\FloatBarrier
\HsweepFig{10.0}{1.0}{1.0}
\UsweepFig{10.0}{1.0}{1.0}
\HsweepFig{10.0}{1.0}{2.0}
\UsweepFig{10.0}{1.0}{2.0}
\HsweepFig{10.0}{1.0}{3.0}
\UsweepFig{10.0}{1.0}{3.0}
\HsweepFig{10.0}{1.0}{4.0}
\UsweepFig{10.0}{1.0}{4.0}

\FloatBarrier
\HsweepFig{5.0}{1.0}{2.0}
\UsweepFig{5.0}{1.0}{2.0}

\end{document}
